\chapter{Binary Trees}
\chapterepigraph{A binary tree is recursion made visible. Almost every
problem in this chapter has the same shape underneath: solve the left
subtree, solve the right subtree, then combine those two answers into an
answer for the current node. Master that single ``ask both children, then
decide'' rhythm and the entire chapter -- traversals, heights, diameters,
views, path sums, reconstruction -- collapses into variations on one
idea.}

\section{The Shape of Every Tree Problem}
\label{sec:trees-intro}

A \textbf{binary tree} is either empty, or a \textbf{node} holding a value
together with two children, each itself the root of a binary tree. That
recursive definition is not just a formality -- it is the template for
almost every algorithm here. A node knows nothing about the tree except
its own value and two pointers; any global fact (the height, the diameter,
whether the tree is balanced, the maximum path sum) must be assembled by
asking each child for its answer and combining.

\begin{ideabox}[The One Rhythm Behind This Whole Chapter]
For a huge fraction of tree problems the skeleton is identical:
\[
  \textit{solve}(node) = \textit{combine}\big(\textit{solve}(node.left),\
  \textit{solve}(node.right),\ node.val\big).
\]
The only thing that changes from problem to problem is \emph{what each
call returns} and \emph{how the two child answers are combined}. Height
returns $1+\max$; ``is balanced'' returns height but also checks a
condition; diameter returns height but also updates a global best;
max-path-sum returns the best downward path but updates a global best
using both sides. Learn to see every problem as ``what do I return, and
how do I combine,'' and the code writes itself.
\end{ideabox}

\subsection*{How to Recognise Which Technique a Tree Problem Needs}

\begin{patternbox}[Recurring Binary-Tree Keywords]
\begin{itemize}
  \item \textbf{``Visit/print every node in some order''} -- a
        \textbf{traversal}. Depth-first (preorder, inorder, postorder)
        is plain recursion, or an explicit stack when an interviewer
        forbids recursion; breadth-first (level order) is a
        \textbf{queue}.
  \item \textbf{``Height, depth, balanced, diameter, maximum path
        sum''} -- \textbf{bottom-up (postorder) recursion}: each node's
        answer needs both children's answers \emph{first}, so compute
        children, then the node. A single global variable often records
        the best answer while each call returns the value its parent
        needs.
  \item \textbf{``Top view, bottom view, vertical order, width''} --
        \textbf{level-order (BFS) with coordinates}: tag each node with a
        horizontal distance (and sometimes depth), then read the map off
        in the required order.
  \item \textbf{``Left view, right view, level order, zigzag''} --
        \textbf{BFS level by level}, picking the first/last node of each
        level or alternating direction.
  \item \textbf{``Lowest common ancestor, root-to-node path, nodes at
        distance K, minimum time to burn''} -- \textbf{path/relationship
        problems}: either a single recursion that reports back up whether
        a target was found, or a preprocessing pass that records each
        node's parent so the tree can be walked in \emph{any} direction.
  \item \textbf{``Construct a tree from its traversals,
        serialize/deserialize''} -- \textbf{reconstruction}: one
        traversal fixes the root, another (or a delimiter scheme) splits
        the remainder into left and right, applied recursively.
  \item \textbf{``$O(1)$ space traversal, no stack, no recursion''} --
        \textbf{Morris traversal}, which borrows the tree's own null
        right pointers as temporary threads back to ancestors.
\end{itemize}
\end{patternbox}

Throughout this chapter a node is the following structure, and every code
listing is a complete, compilable program:

\begin{lstlisting}
struct TreeNode {
    int val;
    TreeNode* left;
    TreeNode* right;
    TreeNode(int v) : val(v), left(nullptr), right(nullptr) {}
};
\end{lstlisting}

Each section follows the established structure: problem statement and
keywords; the reasoning that identifies the right traversal or recursive
technique and why it is correct; the algorithm in words and pseudocode; a
hand-traced dry run; complete compilable C++ code; and a time/space
complexity analysis.

\newpage

% ---- Traversals ----
\section{Recursive Traversals: Preorder, Inorder, Postorder}
\label{sec:trees-rec-traversals}

\problemstatement{Given the root of a binary tree, return its preorder,
inorder, and postorder traversals. Preorder visits \emph{node, left,
right}; inorder visits \emph{left, node, right}; postorder visits
\emph{left, right, node}.}

\begin{patternbox}
\emph{``Visit every node in a fixed relative order''} is the definition of
a depth-first traversal. The only thing separating the three is
\textbf{where the current node is printed relative to its two recursive
calls} -- before both (pre), between them (in), or after both (post). If
you can write one, you can write all three by moving a single line.
\end{patternbox}

\subsection*{Understanding the problem carefully}

The recursive structure of a binary tree hands you the traversal almost
for free. To process a whole tree you process its root and its two
subtrees; each subtree is processed the same way. The word
``pre/in/post'' refers only to \emph{when the root itself is recorded}
relative to recursing into the left and right subtrees.

\begin{ideabox}[One Template, Three Positions]
Every depth-first traversal is the same three actions -- \textsc{record
node}, \textsc{recurse left}, \textsc{recurse right} -- performed in a
different order:
\begin{itemize}
  \item \textbf{Preorder:} record, left, right.
  \item \textbf{Inorder:} left, record, right.
  \item \textbf{Postorder:} left, right, record.
\end{itemize}
The recursion itself is byte-for-byte identical; only the placement of the
``record'' step moves.
\end{ideabox}

\subsection*{Algorithm}

\begin{enumerate}
  \item \textsc{Inorder}(\textit{node}, \textit{out}):
  \begin{enumerate}
    \item If \textit{node} is null, return.
    \item \textsc{Inorder}(\textit{node}.left, \textit{out}).
    \item Append \textit{node}.val to \textit{out}.
    \item \textsc{Inorder}(\textit{node}.right, \textit{out}).
  \end{enumerate}
  \item Preorder and postorder are the same three lines with the append
        moved to the top or the bottom respectively.
\end{enumerate}

\subsection*{Dry run}

\dryrun{Tree: root $1$, right child $2$, and $2$'s left child $3$.}

\begin{itemize}
  \item \textbf{Preorder}: record $1$; left of $1$ is null; go right to
        $2$; record $2$; left of $2$ is $3$; record $3$ (its children
        null); right of $2$ null. Result: $1,2,3$.
  \item \textbf{Inorder}: left of $1$ null; record $1$; go right to $2$;
        left of $2$ is $3$, record $3$; record $2$; right null. Result:
        $1,3,2$.
  \item \textbf{Postorder}: left of $1$ null; right subtree of $1$: left
        of $2$ is $3$, record $3$; right of $2$ null; record $2$; record
        $1$. Result: $3,2,1$.
\end{itemize}

\subsection*{C++ implementation}

\begin{lstlisting}
#include <bits/stdc++.h>
using namespace std;

struct TreeNode {
    int val;
    TreeNode* left;
    TreeNode* right;
    TreeNode(int v) : val(v), left(nullptr), right(nullptr) {}
};

void preorder(TreeNode* node, vector<int>& out) {
    if (node == nullptr) return;
    out.push_back(node->val);   // record
    preorder(node->left, out);
    preorder(node->right, out);
}

void inorder(TreeNode* node, vector<int>& out) {
    if (node == nullptr) return;
    inorder(node->left, out);
    out.push_back(node->val);   // record
    inorder(node->right, out);
}

void postorder(TreeNode* node, vector<int>& out) {
    if (node == nullptr) return;
    postorder(node->left, out);
    postorder(node->right, out);
    out.push_back(node->val);   // record
}

int main() {
    TreeNode* root = new TreeNode(1);
    root->right = new TreeNode(2);
    root->right->left = new TreeNode(3);

    vector<int> pre, in, post;
    preorder(root, pre);
    inorder(root, in);
    postorder(root, post);

    auto printv = [](const string& name, vector<int>& v) {
        cout << name << ": ";
        for (int x : v) cout << x << " ";
        cout << "\n";
    };
    printv("preorder ", pre);
    printv("inorder  ", in);
    printv("postorder", post);
    return 0;
}
\end{lstlisting}

\begin{complexitybox}
\textbf{Time Complexity.} $O(n)$ -- each node is visited exactly once.
\textbf{Space Complexity.} $O(h)$ for the recursion stack, where $h$ is
the tree's height ($O(n)$ in the worst case of a skewed tree, $O(\log n)$
when balanced), plus $O(n)$ for the output array.
\end{complexitybox}

\begin{takeawaybox}
The three depth-first traversals are one algorithm with the ``record''
step in three different positions. Every later depth-first problem in this
chapter is a traversal with extra work grafted onto that record step, so
internalise this skeleton before moving on.
\end{takeawaybox}

\section{Level Order Traversal (BFS)}
\label{sec:trees-level-order}

\problemstatement{Given the root of a binary tree, return its level order
traversal: the node values grouped level by level, from the root's level
downward, left to right within each level.}

\begin{patternbox}
\emph{``Level by level,'' ``breadth-first,'' ``top-to-bottom''} are the
signature of a \textbf{queue}-based traversal. Whenever a tree problem
asks for anything organised by \emph{depth} -- level order, left/right
view, bottom/top view, maximum width, minimum time to burn -- a BFS with a
queue is almost always the backbone.
\end{patternbox}

\subsection*{Understanding the problem carefully}

Depth-first recursion dives all the way down one branch before touching
the next, so it visits nodes out of level order. To visit nodes in order
of increasing depth we need the opposite discipline: fully process one
level before starting the next. A \textbf{FIFO queue} does exactly this --
if we always enqueue a node's children when we dequeue it, then all of
level $k$ is enqueued before any node of level $k+1$, so the queue hands
nodes back in level order automatically.

\begin{ideabox}[Freeze the Level Size Before Draining It]
When you begin processing a level, the queue holds \emph{exactly} the
nodes of that level and nothing else. Record its size $s$ first, then
dequeue precisely $s$ nodes: those $s$ are the current level, and the
children you enqueue while draining them form the next level, cleanly
separated. This ``read the size, then loop that many times'' trick is what
lets a single queue produce grouped-by-level output.
\end{ideabox}

\subsection*{Algorithm}

\begin{enumerate}
  \item If root is null, return empty.
  \item Push root into a queue.
  \item While the queue is non-empty:
  \begin{enumerate}
    \item Let $s$ be the current queue size; start a new level list.
    \item Repeat $s$ times: pop the front node, append its value to the
          level list, and push its non-null children.
    \item Append the level list to the answer.
  \end{enumerate}
\end{enumerate}

\subsection*{Dry run}

\dryrun{Tree: root $1$; children $2$ (left) and $3$ (right); $2$'s
children $4,5$.}

\begin{itemize}
  \item Queue $=[1]$, size $1$: pop $1$, push $2,3$. Level $=[1]$.
  \item Queue $=[2,3]$, size $2$: pop $2$ (push $4,5$), pop $3$. Level
        $=[2,3]$.
  \item Queue $=[4,5]$, size $2$: pop $4$, pop $5$. Level $=[4,5]$.
  \item Result: $[[1],[2,3],[4,5]]$.
\end{itemize}

\subsection*{C++ implementation}

\begin{lstlisting}
#include <bits/stdc++.h>
using namespace std;

struct TreeNode {
    int val;
    TreeNode* left;
    TreeNode* right;
    TreeNode(int v) : val(v), left(nullptr), right(nullptr) {}
};

vector<vector<int>> levelOrder(TreeNode* root) {
    vector<vector<int>> ans;
    if (root == nullptr) return ans;

    queue<TreeNode*> q;
    q.push(root);
    while (!q.empty()) {
        int s = q.size();
        vector<int> level;
        for (int i = 0; i < s; i++) {
            TreeNode* node = q.front();
            q.pop();
            level.push_back(node->val);
            if (node->left)  q.push(node->left);
            if (node->right) q.push(node->right);
        }
        ans.push_back(level);
    }
    return ans;
}

int main() {
    TreeNode* root = new TreeNode(1);
    root->left = new TreeNode(2);
    root->right = new TreeNode(3);
    root->left->left = new TreeNode(4);
    root->left->right = new TreeNode(5);

    for (auto& level : levelOrder(root)) {
        for (int x : level) cout << x << " ";
        cout << "\n";
    }
    return 0;
}
\end{lstlisting}

\begin{complexitybox}
\textbf{Time Complexity.} $O(n)$ -- every node is enqueued and dequeued
exactly once. \textbf{Space Complexity.} $O(n)$ for the queue, which in
the worst case (a full last level) holds up to $n/2$ nodes at once.
\end{complexitybox}

\begin{takeawaybox}
The queue plus the ``freeze the level size'' loop is the single reusable
engine behind every by-depth view in this chapter. Get it fluent here;
top view, bottom view, right view and maximum width all reuse it verbatim
with a small tweak to what gets recorded.
\end{takeawaybox}

\section{Iterative Preorder Traversal}
\label{sec:trees-iter-preorder}

\problemstatement{Return the preorder traversal (node, left, right) of a
binary tree \emph{without} recursion, using an explicit stack.}

\begin{patternbox}
\emph{``Do this traversal iteratively / without recursion''} means:
simulate the call stack yourself with an explicit \textbf{stack}. Preorder
is the easiest of the three to convert because the node is recorded the
instant it is seen, before either child -- so there is no ``come back to
this node later'' bookkeeping to manage.
\end{patternbox}

\subsection*{Understanding the problem carefully}

Recursion works by an implicit stack of pending calls. To remove the
recursion we make that stack explicit. In preorder we want to output a
node, then fully handle its left subtree, then its right subtree. A stack
(LIFO) lets us achieve exactly this if we are careful about push order.

\begin{ideabox}[Push Right Before Left So Left Pops First]
Pop a node and immediately record it. Then push its \emph{right} child and
\emph{then} its \emph{left} child. Because a stack returns the most
recently pushed item first, the left child is processed before the right
-- matching preorder's ``node, then all of left, then all of right.''
Pushing right first is the whole trick; reverse it and you would get
``node, right, left.''
\end{ideabox}

\subsection*{Algorithm}

\begin{enumerate}
  \item If root is null, return empty. Push root onto a stack.
  \item While the stack is non-empty:
  \begin{enumerate}
    \item Pop a node; append its value to the output.
    \item If it has a right child, push it.
    \item If it has a left child, push it.
  \end{enumerate}
\end{enumerate}

\subsection*{Dry run}

\dryrun{Tree: root $1$; children $2$ (left), $3$ (right); $2$'s left child
$4$.}

\begin{itemize}
  \item Stack $=[1]$: pop $1$, record $1$, push $3$ then $2$. Stack
        $=[3,2]$ (top is $2$).
  \item Pop $2$, record $2$, push (no right), push $4$. Stack $=[3,4]$.
  \item Pop $4$, record $4$ (no children). Stack $=[3]$.
  \item Pop $3$, record $3$. Stack empty.
  \item Result: $1,2,4,3$.
\end{itemize}

\subsection*{C++ implementation}

\begin{lstlisting}
#include <bits/stdc++.h>
using namespace std;

struct TreeNode {
    int val;
    TreeNode* left;
    TreeNode* right;
    TreeNode(int v) : val(v), left(nullptr), right(nullptr) {}
};

vector<int> preorderIterative(TreeNode* root) {
    vector<int> out;
    if (root == nullptr) return out;

    stack<TreeNode*> st;
    st.push(root);
    while (!st.empty()) {
        TreeNode* node = st.top();
        st.pop();
        out.push_back(node->val);
        if (node->right) st.push(node->right);   // right first
        if (node->left)  st.push(node->left);    // so left pops first
    }
    return out;
}

int main() {
    TreeNode* root = new TreeNode(1);
    root->left = new TreeNode(2);
    root->right = new TreeNode(3);
    root->left->left = new TreeNode(4);

    for (int x : preorderIterative(root)) cout << x << " ";
    cout << "\n";
    return 0;
}
\end{lstlisting}

\begin{complexitybox}
\textbf{Time Complexity.} $O(n)$ -- each node is pushed and popped exactly
once. \textbf{Space Complexity.} $O(h)$ for the stack in the worst case
(the stack never holds more than one node per level of the current path
plus pending right children).
\end{complexitybox}

\begin{takeawaybox}
Preorder is the friendliest traversal to make iterative because ``record''
happens on first sight. The push-right-then-left order is the reusable
detail. Inorder and postorder are harder precisely because their record
step is delayed until after one or both subtrees are done.
\end{takeawaybox}

\section{Iterative Inorder Traversal}
\label{sec:trees-iter-inorder}

\problemstatement{Return the inorder traversal (left, node, right) of a
binary tree without recursion, using an explicit stack.}

\begin{patternbox}
Inorder's record step is \emph{delayed}: a node cannot be output until its
entire left subtree is done. So the iterative version has two phases that
alternate -- \textbf{dive left pushing everything}, then \textbf{pop,
record, and turn right}. This ``go as far left as possible, then process,
then step right'' loop is worth memorising; it reappears in the BST
iterator.
\end{patternbox}

\subsection*{Understanding the problem carefully}

In inorder we must reach the leftmost node before recording anything. The
stack stores the chain of ancestors we still owe a visit to. We push
nodes as we descend left; when we can descend no further, the top of the
stack is the next node to record; after recording it we move to its right
subtree and repeat the leftward dive from there.

\begin{ideabox}[A Pointer and a Stack, Cooperating]
Keep a running pointer \textit{curr}. While \textit{curr} is non-null,
push it and move \textit{curr} to its left child (the leftward dive). When
\textit{curr} becomes null, pop a node, record it, and set \textit{curr}
to that node's right child (turn right, then dive left again). The loop
ends only when both the stack is empty \emph{and} \textit{curr} is null.
\end{ideabox}

\subsection*{Algorithm}

\begin{enumerate}
  \item Set \textit{curr} $=$ root; use an empty stack.
  \item While \textit{curr} is non-null \textbf{or} the stack is
        non-empty:
  \begin{enumerate}
    \item While \textit{curr} is non-null: push \textit{curr}, set
          \textit{curr} $=$ \textit{curr}.left.
    \item Pop a node; append its value; set \textit{curr} $=$
          node.right.
  \end{enumerate}
\end{enumerate}

\subsection*{Dry run}

\dryrun{Tree: root $2$; left child $1$; right child $3$.}

\begin{itemize}
  \item \textit{curr}$=2$: push $2$, \textit{curr}$=1$. Push $1$,
        \textit{curr}$=$null.
  \item Pop $1$, record $1$, \textit{curr}$=$null (right of $1$).
  \item Pop $2$, record $2$, \textit{curr}$=3$.
  \item Push $3$, \textit{curr}$=$null. Pop $3$, record $3$.
  \item Result: $1,2,3$.
\end{itemize}

\subsection*{C++ implementation}

\begin{lstlisting}
#include <bits/stdc++.h>
using namespace std;

struct TreeNode {
    int val;
    TreeNode* left;
    TreeNode* right;
    TreeNode(int v) : val(v), left(nullptr), right(nullptr) {}
};

vector<int> inorderIterative(TreeNode* root) {
    vector<int> out;
    stack<TreeNode*> st;
    TreeNode* curr = root;

    while (curr != nullptr || !st.empty()) {
        while (curr != nullptr) {   // dive left
            st.push(curr);
            curr = curr->left;
        }
        curr = st.top();            // leftmost pending node
        st.pop();
        out.push_back(curr->val);   // record
        curr = curr->right;         // turn right
    }
    return out;
}

int main() {
    TreeNode* root = new TreeNode(2);
    root->left = new TreeNode(1);
    root->right = new TreeNode(3);

    for (int x : inorderIterative(root)) cout << x << " ";
    cout << "\n";
    return 0;
}
\end{lstlisting}

\begin{complexitybox}
\textbf{Time Complexity.} $O(n)$ -- each node is pushed once and popped
once. \textbf{Space Complexity.} $O(h)$ -- the stack holds at most one
node per level along the current leftmost chain.
\end{complexitybox}

\begin{takeawaybox}
The ``dive left, pop-record, turn right'' loop is the canonical iterative
inorder and the exact engine behind a BST iterator that returns sorted
values one at a time. Memorise the two-part loop shape; the same structure
handles ``next smallest'' queries on demand.
\end{takeawaybox}

\section{Iterative Postorder Traversal (Two Stacks)}
\label{sec:trees-iter-post-two}

\problemstatement{Return the postorder traversal (left, right, node) of a
binary tree without recursion, using two stacks.}

\begin{patternbox}
Postorder is the hardest to make iterative, because a node is recorded
only \emph{after both} subtrees. The two-stack method sidesteps the
difficulty with a slick observation: postorder is the exact
\textbf{reverse} of a slightly modified preorder. If you can produce
``node, right, left'' and then reverse it, you have postorder for free.
\end{patternbox}

\subsection*{Understanding the problem carefully}

Postorder is left, right, node. Read it backwards and you get node, right,
left -- which is just preorder with the two child visits swapped. So if we
run a preorder-like traversal that pushes \emph{left before right} (giving
node, right, left order) and collect the results, then reverse the whole
collection, we recover true postorder. The second stack is simply the
container we reverse.

\begin{ideabox}[Postorder = reverse of (node, right, left)]
Run a preorder variant using stack~1 that emits nodes in \emph{node,
right, left} order, pushing each emitted node onto stack~2. Popping
stack~2 at the end yields those nodes in reverse -- i.e.\ \emph{left,
right, node}, exactly postorder. No delayed-visit bookkeeping is needed.
\end{ideabox}

\subsection*{Algorithm}

\begin{enumerate}
  \item If root is null, return empty. Push root onto stack~1.
  \item While stack~1 is non-empty:
  \begin{enumerate}
    \item Pop a node from stack~1 and push it onto stack~2.
    \item Push that node's \emph{left} child, then its \emph{right}
          child, onto stack~1 (if they exist).
  \end{enumerate}
  \item Pop everything off stack~2 into the output; that is the
        postorder.
\end{enumerate}

\subsection*{Dry run}

\dryrun{Tree: root $1$; children $2$ (left), $3$ (right).}

\begin{itemize}
  \item Stack1 $=[1]$: pop $1\to$ stack2 $=[1]$; push $2$ then $3$. Stack1
        $=[2,3]$.
  \item Pop $3\to$ stack2 $=[1,3]$ (no children). Stack1 $=[2]$.
  \item Pop $2\to$ stack2 $=[1,3,2]$. Stack1 empty.
  \item Empty stack2: $2,3,1$. That is the postorder.
\end{itemize}

\subsection*{C++ implementation}

\begin{lstlisting}
#include <bits/stdc++.h>
using namespace std;

struct TreeNode {
    int val;
    TreeNode* left;
    TreeNode* right;
    TreeNode(int v) : val(v), left(nullptr), right(nullptr) {}
};

vector<int> postorderTwoStacks(TreeNode* root) {
    vector<int> out;
    if (root == nullptr) return out;

    stack<TreeNode*> s1, s2;
    s1.push(root);
    while (!s1.empty()) {
        TreeNode* node = s1.top();
        s1.pop();
        s2.push(node);
        if (node->left)  s1.push(node->left);
        if (node->right) s1.push(node->right);
    }
    while (!s2.empty()) {
        out.push_back(s2.top()->val);
        s2.pop();
    }
    return out;
}

int main() {
    TreeNode* root = new TreeNode(1);
    root->left = new TreeNode(2);
    root->right = new TreeNode(3);

    for (int x : postorderTwoStacks(root)) cout << x << " ";
    cout << "\n";
    return 0;
}
\end{lstlisting}

\begin{complexitybox}
\textbf{Time Complexity.} $O(n)$ -- every node passes through both stacks
once. \textbf{Space Complexity.} $O(n)$ -- stack~2 accumulates all $n$
nodes before the final reversal, which is the price of the simpler logic.
\end{complexitybox}

\begin{takeawaybox}
Recognising that postorder is a reversed, child-swapped preorder turns a
fiddly problem into two easy passes. The trade is memory: it holds all $n$
nodes at once. When an interviewer demands $O(h)$ space, use the one-stack
method in the next section instead.
\end{takeawaybox}

\section{Iterative Postorder Traversal (One Stack)}
\label{sec:trees-iter-post-one}

\problemstatement{Return the postorder traversal of a binary tree without
recursion, using a single stack and $O(h)$ auxiliary space.}

\begin{patternbox}
The tell-tale that a one-stack postorder is wanted is a space constraint:
the two-stack method is simpler but stores all $n$ nodes. With one stack we
must handle the genuine difficulty head-on -- \emph{a node may sit on the
stack while we visit first its left subtree and only later its right} --
so we need a way to know, when a node resurfaces, whether we are returning
from its left or its right side.
\end{patternbox}

\subsection*{Understanding the problem carefully}

In postorder we record a node only after \emph{both} its children are
done. When we pop back up to a node, we must distinguish ``I have finished
its left subtree, now go right'' from ``I have finished its right subtree,
now record it.'' Tracking the \textbf{last node we recorded} resolves this
ambiguity: if that last-recorded node is the current node's right child
(or the right child is null), both subtrees are complete and the node may
be recorded; otherwise there is still a right subtree to descend into.

\begin{ideabox}[Remember the Last Node You Output]
Keep \textit{curr} (the cursor) and \textit{lastVisited} (the most
recently recorded node). Descend left, pushing. At the top of the stack:
if its right child exists and has not just been visited, move right;
otherwise record the top, pop it, and set \textit{lastVisited} to it. That
single comparison against \textit{lastVisited} tells you which side you are
returning from.
\end{ideabox}

\subsection*{Algorithm}

\begin{enumerate}
  \item \textit{curr} $=$ root; \textit{lastVisited} $=$ null; empty
        stack.
  \item While \textit{curr} is non-null or the stack is non-empty:
  \begin{enumerate}
    \item While \textit{curr} is non-null: push \textit{curr},
          \textit{curr} $=$ \textit{curr}.left.
    \item Let \textit{peek} $=$ top of stack. If \textit{peek}.right
          exists and \textit{peek}.right $\ne$ \textit{lastVisited}: set
          \textit{curr} $=$ \textit{peek}.right.
    \item Else: record \textit{peek}, pop it, set \textit{lastVisited}
          $=$ \textit{peek}.
  \end{enumerate}
\end{enumerate}

\subsection*{Dry run}

\dryrun{Tree: root $1$; left $2$; right $3$.}

\begin{itemize}
  \item Dive left from $1$: push $1$, push $2$. \textit{curr}$=$null.
  \item Peek $2$: right is null $\Rightarrow$ record $2$, pop,
        \textit{lastVisited}$=2$.
  \item Peek $1$: right $3\ne$\textit{lastVisited} $\Rightarrow$
        \textit{curr}$=3$. Push $3$.
  \item Peek $3$: right null $\Rightarrow$ record $3$, pop,
        \textit{lastVisited}$=3$.
  \item Peek $1$: right $3=$\textit{lastVisited} $\Rightarrow$ record $1$,
        pop.
  \item Result: $2,3,1$.
\end{itemize}

\subsection*{C++ implementation}

\begin{lstlisting}
#include <bits/stdc++.h>
using namespace std;

struct TreeNode {
    int val;
    TreeNode* left;
    TreeNode* right;
    TreeNode(int v) : val(v), left(nullptr), right(nullptr) {}
};

vector<int> postorderOneStack(TreeNode* root) {
    vector<int> out;
    stack<TreeNode*> st;
    TreeNode* curr = root;
    TreeNode* lastVisited = nullptr;

    while (curr != nullptr || !st.empty()) {
        while (curr != nullptr) {
            st.push(curr);
            curr = curr->left;
        }
        TreeNode* peek = st.top();
        if (peek->right != nullptr && peek->right != lastVisited) {
            curr = peek->right;             // still owe the right subtree
        } else {
            out.push_back(peek->val);       // both sides done: record
            lastVisited = peek;
            st.pop();
        }
    }
    return out;
}

int main() {
    TreeNode* root = new TreeNode(1);
    root->left = new TreeNode(2);
    root->right = new TreeNode(3);

    for (int x : postorderOneStack(root)) cout << x << " ";
    cout << "\n";
    return 0;
}
\end{lstlisting}

\begin{complexitybox}
\textbf{Time Complexity.} $O(n)$ -- each node is pushed once, and the
peek/right check amortises to $O(1)$ per node. \textbf{Space Complexity.}
$O(h)$ -- the stack only ever holds the current root-to-node path, unlike
the two-stack version's $O(n)$.
\end{complexitybox}

\begin{takeawaybox}
When you must return from a node's left versus right side and act
differently, a \textit{lastVisited} pointer is the standard way to tell
them apart without extra per-node flags. This ``remember where you came
from'' idea generalises to any iterative post-processing traversal.
\end{takeawaybox}

\section{Preorder, Inorder, and Postorder in One Traversal}
\label{sec:trees-all-three}

\problemstatement{Compute all three depth-first traversals -- preorder,
inorder, and postorder -- in a \emph{single} pass over the tree, using one
explicit stack.}

\begin{patternbox}
\emph{``All three in one traversal''} is a state-machine problem. Each node
is genuinely visited three times during a depth-first walk: on the way
down (before its left subtree), between its subtrees, and on the way up
(after its right subtree). Attaching a \textbf{state counter $1/2/3$} to
each stacked node lets one loop record the node into the right list at the
right moment.
\end{patternbox}

\subsection*{Understanding the problem carefully}

Picture the classic ``draw a line around the tree'' visualisation of a
DFS: your pen touches each node three times -- once on its left, once
underneath, once on its right. Preorder records the first touch, inorder
the second, postorder the third. If we stack each node together with a
counter saying which touch is next, a single loop can dispatch it to the
correct output list and advance its counter.

\begin{ideabox}[One Node, Three Touches, Three States]
Push each node with state $1$. When popped:
\begin{itemize}
  \item \textbf{State 1} (first touch): add to \emph{preorder}; increment
        its state to $2$ and re-push it; then push its left child (state
        $1$).
  \item \textbf{State 2} (second touch): add to \emph{inorder}; increment
        to $3$ and re-push; then push its right child (state $1$).
  \item \textbf{State 3} (third touch): add to \emph{postorder}; discard
        the node.
\end{itemize}
\end{ideabox}

\subsection*{Algorithm}

\begin{enumerate}
  \item Push $(\text{root}, 1)$ onto a stack. While non-empty, pop
        $(node, state)$:
  \begin{enumerate}
    \item If $state=1$: append to preorder; push $(node,2)$; if left
          child exists push $(\text{left},1)$.
    \item If $state=2$: append to inorder; push $(node,3)$; if right
          child exists push $(\text{right},1)$.
    \item If $state=3$: append to postorder.
  \end{enumerate}
\end{enumerate}

\subsection*{Dry run}

\dryrun{Tree: root $1$; left $2$; right $3$.}

\begin{itemize}
  \item Pop $(1,1)$: pre$=[1]$; push $(1,2)$; push $(2,1)$.
  \item Pop $(2,1)$: pre$=[1,2]$; push $(2,2)$ (no left).
  \item Pop $(2,2)$: in$=[2]$; push $(2,3)$ (no right).
  \item Pop $(2,3)$: post$=[2]$.
  \item Pop $(1,2)$: in$=[2,1]$; push $(1,3)$; push $(3,1)$.
  \item Pop $(3,1)$: pre$=[1,2,3]$; push $(3,2)$. Pop $(3,2)$:
        in$=[2,1,3]$; push $(3,3)$. Pop $(3,3)$: post$=[2,3]$.
  \item Pop $(1,3)$: post$=[2,3,1]$.
  \item Pre $=1,2,3$; \; In $=2,1,3$; \; Post $=2,3,1$.
\end{itemize}

\subsection*{C++ implementation}

\begin{lstlisting}
#include <bits/stdc++.h>
using namespace std;

struct TreeNode {
    int val;
    TreeNode* left;
    TreeNode* right;
    TreeNode(int v) : val(v), left(nullptr), right(nullptr) {}
};

void allThree(TreeNode* root, vector<int>& pre,
              vector<int>& in, vector<int>& post) {
    if (root == nullptr) return;
    stack<pair<TreeNode*, int>> st;
    st.push({root, 1});

    while (!st.empty()) {
        auto [node, state] = st.top();
        st.pop();
        if (state == 1) {
            pre.push_back(node->val);
            st.push({node, 2});
            if (node->left) st.push({node->left, 1});
        } else if (state == 2) {
            in.push_back(node->val);
            st.push({node, 3});
            if (node->right) st.push({node->right, 1});
        } else {
            post.push_back(node->val);
        }
    }
}

int main() {
    TreeNode* root = new TreeNode(1);
    root->left = new TreeNode(2);
    root->right = new TreeNode(3);

    vector<int> pre, in, post;
    allThree(root, pre, in, post);

    auto printv = [](const string& n, vector<int>& v) {
        cout << n << ": ";
        for (int x : v) cout << x << " ";
        cout << "\n";
    };
    printv("pre ", pre);
    printv("in  ", in);
    printv("post", post);
    return 0;
}
\end{lstlisting}

\begin{complexitybox}
\textbf{Time Complexity.} $O(n)$ -- each node is pushed and popped exactly
three times, a constant factor. \textbf{Space Complexity.} $O(h)$ for the
stack, plus $O(n)$ for the three output arrays.
\end{complexitybox}

\begin{takeawaybox}
The state $1/2/3$ per stack frame is exactly what a recursive DFS stores
implicitly as ``which line am I on when I return.'' Making that state
explicit is a general technique for merging several traversal-order tasks
into one pass.
\end{takeawaybox}


% ---- Medium ----
\section{Maximum Depth (Height) of a Binary Tree}
\label{sec:trees-max-depth}

\problemstatement{Given the root of a binary tree, return its maximum
depth: the number of nodes along the longest path from the root down to a
leaf.}

\begin{patternbox}
\emph{``Height,'' ``depth,'' ``longest path to a leaf''} is the archetypal
\textbf{bottom-up} recursion, and the first concrete instance of this
chapter's core rhythm. A node's height depends on its children's heights,
so you must have both child answers \emph{before} you can answer for the
node -- that ordering \emph{is} postorder.
\end{patternbox}

\subsection*{Understanding the problem carefully}

The height of a tree rooted at a node is $1$ (for the node itself) plus
the height of its taller subtree. An empty tree has height $0$. This is
the purest example of ``ask both children, combine'': the combine step is
simply $1+\max(\text{left height}, \text{right height})$.

\begin{ideabox}[Height = 1 + max of the Two Child Heights]
Recurse to the left and right children, obtain their heights, and return
one more than the larger. The base case (null returns $0$) is what
terminates the recursion and seeds the count. No global state is needed:
the answer flows straight up the return values.
\end{ideabox}

\subsection*{Algorithm}

\begin{enumerate}
  \item \textsc{Height}(\textit{node}):
  \begin{enumerate}
    \item If \textit{node} is null, return $0$.
    \item $L = $ \textsc{Height}(\textit{node}.left); $R = $
          \textsc{Height}(\textit{node}.right).
    \item Return $1 + \max(L, R)$.
  \end{enumerate}
\end{enumerate}

\subsection*{Dry run}

\dryrun{Tree: root $1$; left $2$ (leaf); right $3$ whose right child is
$4$ (leaf).}

\begin{itemize}
  \item Height($2$)$=1$ (both children null $\to 1+\max(0,0)$).
  \item Height($4$)$=1$. Height($3$)$=1+\max(0,1)=2$.
  \item Height($1$)$=1+\max(1,2)=3$.
\end{itemize}

\subsection*{C++ implementation}

\begin{lstlisting}
#include <bits/stdc++.h>
using namespace std;

struct TreeNode {
    int val;
    TreeNode* left;
    TreeNode* right;
    TreeNode(int v) : val(v), left(nullptr), right(nullptr) {}
};

int maxDepth(TreeNode* node) {
    if (node == nullptr) return 0;
    int L = maxDepth(node->left);
    int R = maxDepth(node->right);
    return 1 + max(L, R);
}

int main() {
    TreeNode* root = new TreeNode(1);
    root->left = new TreeNode(2);
    root->right = new TreeNode(3);
    root->right->right = new TreeNode(4);

    cout << "Max depth: " << maxDepth(root) << "\n";
    return 0;
}
\end{lstlisting}

\begin{complexitybox}
\textbf{Time Complexity.} $O(n)$ -- each node contributes $O(1)$ combine
work. \textbf{Space Complexity.} $O(h)$ for the recursion stack.
\end{complexitybox}

\begin{takeawaybox}
Height is the ``hello world'' of bottom-up tree recursion, and the exact
subroutine reused inside balanced-tree, diameter, and max-path-sum
problems. Each of those is ``compute height, but also check or update
something extra along the way.''
\end{takeawaybox}

\section{Check for a Balanced Binary Tree}
\label{sec:trees-balanced}

\problemstatement{A binary tree is height-balanced if, for \emph{every}
node, the heights of its left and right subtrees differ by at most $1$.
Given the root, determine whether the tree is balanced.}

\begin{patternbox}
The naive reading -- ``for every node, compute both subtree heights and
compare'' -- is an $O(n^2)$ trap on a skewed tree, because height is
recomputed from scratch at every node. The keyword \emph{``for every
node''} combined with \emph{``heights''} says: compute height once,
bottom-up, and \textbf{fold the balance check into that same pass}.
\end{patternbox}

\subsection*{Understanding the problem carefully}

Height and the balance condition are computed from the same recursion. If
each recursive call returns the height of its subtree, the parent can, in
$O(1)$, check whether its two child heights differ by more than $1$ before
returning its own height. A single sentinel value (say $-1$) propagated
upward can signal ``a violation was found below'' so the whole recursion
short-circuits.

\begin{ideabox}[Return Height, or $-1$ to Mean ``Already Unbalanced'']
Write a helper that returns the subtree height normally, but returns $-1$
the moment it detects imbalance -- either because a child already returned
$-1$, or because the two child heights differ by more than $1$. The final
answer is ``balanced'' iff the root's call returns something other than
$-1$. One pass, $O(n)$.
\end{ideabox}

\subsection*{Algorithm}

\begin{enumerate}
  \item \textsc{Check}(\textit{node}):
  \begin{enumerate}
    \item If \textit{node} is null, return $0$.
    \item $L = $ \textsc{Check}(\textit{node}.left); if $L=-1$ return
          $-1$.
    \item $R = $ \textsc{Check}(\textit{node}.right); if $R=-1$ return
          $-1$.
    \item If $|L-R|>1$ return $-1$; else return $1+\max(L,R)$.
  \end{enumerate}
  \item The tree is balanced iff \textsc{Check}(root)$\ne -1$.
\end{enumerate}

\subsection*{Dry run}

\dryrun{Tree: root $1$; left $2$ with a left child $3$ that has a left
child $4$; right subtree of $1$ empty.}

\begin{itemize}
  \item Height($4$)$=1$, Height($3$)$=1+\max(1,0)=2$.
  \item At $2$: left height $2$, right height $0$; $|2-0|=2>1
        \Rightarrow$ return $-1$.
  \item $-1$ propagates to $1$; answer: \textbf{not balanced}.
\end{itemize}

\subsection*{C++ implementation}

\begin{lstlisting}
#include <bits/stdc++.h>
using namespace std;

struct TreeNode {
    int val;
    TreeNode* left;
    TreeNode* right;
    TreeNode(int v) : val(v), left(nullptr), right(nullptr) {}
};

int check(TreeNode* node) {
    if (node == nullptr) return 0;
    int L = check(node->left);
    if (L == -1) return -1;
    int R = check(node->right);
    if (R == -1) return -1;
    if (abs(L - R) > 1) return -1;
    return 1 + max(L, R);
}

bool isBalanced(TreeNode* root) {
    return check(root) != -1;
}

int main() {
    TreeNode* root = new TreeNode(1);
    root->left = new TreeNode(2);
    root->left->left = new TreeNode(3);
    root->left->left->left = new TreeNode(4);

    cout << (isBalanced(root) ? "balanced" : "not balanced") << "\n";
    return 0;
}
\end{lstlisting}

\begin{complexitybox}
\textbf{Time Complexity.} $O(n)$ -- height and the balance test share one
postorder pass, versus $O(n^2)$ if height were recomputed per node.
\textbf{Space Complexity.} $O(h)$ for the recursion stack.
\end{complexitybox}

\begin{takeawaybox}
Overloading a return value with a sentinel ($-1$ = ``violation found'') is
a clean way to combine ``compute a quantity'' with ``check a property''
in one traversal, avoiding the quadratic blow-up of recomputing the
quantity at every node.
\end{takeawaybox}

\section{Diameter of a Binary Tree}
\label{sec:trees-diameter}

\problemstatement{The diameter of a binary tree is the length (number of
edges) of the longest path between any two nodes. The path need not pass
through the root. Given the root, return the diameter.}

\begin{patternbox}
\emph{``Longest path between any two nodes, not necessarily through the
root''} is the classic ``return one thing, track another'' pattern. At
each node the \emph{best path passing through it} is left height $+$ right
height, but what the node must \emph{return} to its parent is only its own
height. Two different quantities: a \textbf{global maximum} for the answer,
a \textbf{return value} for the parent.
\end{patternbox}

\subsection*{Understanding the problem carefully}

Any path has a single highest node -- its ``peak.'' The path through a
given peak descends into the left subtree and the right subtree, so its
length is (height of left subtree) $+$ (height of right subtree) edges.
The overall diameter is the maximum of that quantity over all possible
peaks. Since every node is the peak of exactly one candidate path, we
evaluate this candidate at every node during a height computation.

\begin{ideabox}[Compute Height, Update a Global Diameter Along the Way]
Run the ordinary height recursion. At each node, before returning, update
a global answer with $L+R$ (the through-this-node path length in edges).
Then return $1+\max(L,R)$ to the parent as usual. The node contributes its
own candidate to the global max but hands the parent only what the parent
needs.
\end{ideabox}

\subsection*{Algorithm}

\begin{enumerate}
  \item Keep a global \textit{best} $=0$.
  \item \textsc{Height}(\textit{node}):
  \begin{enumerate}
    \item If null, return $0$.
    \item $L=$\textsc{Height}(left), $R=$\textsc{Height}(right).
    \item \textit{best} $=\max(\textit{best},\ L+R)$.
    \item Return $1+\max(L,R)$.
  \end{enumerate}
  \item Answer is \textit{best}.
\end{enumerate}

\subsection*{Dry run}

\dryrun{Tree: root $1$; left $2$ with children $4,5$; right $3$.}

\begin{itemize}
  \item Height($4$)$=$Height($5$)$=1$. At $2$: $L=1,R=1$, best$=\max(0,2)=2$;
        return $2$.
  \item Height($3$)$=1$. At $1$: $L=2,R=1$, best$=\max(2,3)=3$; return
        $3$.
  \item Diameter $=3$ edges (path $4\!-\!2\!-\!1\!-\!3$).
\end{itemize}

\subsection*{C++ implementation}

\begin{lstlisting}
#include <bits/stdc++.h>
using namespace std;

struct TreeNode {
    int val;
    TreeNode* left;
    TreeNode* right;
    TreeNode(int v) : val(v), left(nullptr), right(nullptr) {}
};

int best;

int height(TreeNode* node) {
    if (node == nullptr) return 0;
    int L = height(node->left);
    int R = height(node->right);
    best = max(best, L + R);      // path through this node, in edges
    return 1 + max(L, R);
}

int diameterOfBinaryTree(TreeNode* root) {
    best = 0;
    height(root);
    return best;
}

int main() {
    TreeNode* root = new TreeNode(1);
    root->left = new TreeNode(2);
    root->right = new TreeNode(3);
    root->left->left = new TreeNode(4);
    root->left->right = new TreeNode(5);

    cout << "Diameter: " << diameterOfBinaryTree(root) << "\n";
    return 0;
}
\end{lstlisting}

\begin{complexitybox}
\textbf{Time Complexity.} $O(n)$ -- one height pass with $O(1)$ extra work
per node. A naive ``compute height twice at every node'' solution would be
$O(n^2)$. \textbf{Space Complexity.} $O(h)$ for the recursion stack.
\end{complexitybox}

\begin{takeawaybox}
``Return one value to the parent, update a separate global for the answer''
is the reusable shape for any tree problem where the optimum can bend
through a node rather than continue straight down. The very next section,
maximum path sum, is the same idea with values instead of heights.
\end{takeawaybox}

\section{Maximum Path Sum}
\label{sec:trees-max-path-sum}

\problemstatement{A path is any sequence of nodes connected by edges, where
each adjacent pair shares an edge and no node repeats; it need not pass
through the root. The path sum is the total of its node values. Given the
root (values may be negative), return the maximum path sum.}

\begin{patternbox}
This is the diameter problem with \emph{values instead of edge counts},
plus one subtlety: because values can be negative, a subtree that would
\emph{lower} the total should be dropped entirely. The keywords
\emph{``maximum,'' ``path,'' ``may be negative''} say: bottom-up, track a
global best through each node, but clamp negative child contributions to
zero.
\end{patternbox}

\subsection*{Understanding the problem carefully}

At any node, the best path with that node as its peak is: the node's value,
plus the best downward path into the left subtree (if positive), plus the
best downward path into the right subtree (if positive). That is the
candidate we compare against a global best. But the value we \emph{return}
to the parent can only continue \emph{one} way -- a path cannot fork at the
parent -- so we return the node's value plus the better of its two
downward gains.

\begin{ideabox}[Return a Straight-Down Gain, Score a Bent Path]
Let \textsc{Gain}(node) return the maximum sum of a path that starts at
\textit{node} and goes strictly downward. Compute
$l=\max(0,\textsc{Gain}(\text{left}))$ and
$r=\max(0,\textsc{Gain}(\text{right}))$ (the $\max(0,\cdot)$ drops
harmful subtrees). Update the global best with $node.val+l+r$ (the path
bending through this node). Return $node.val+\max(l,r)$ to the parent (only
one branch may continue).
\end{ideabox}

\subsection*{Algorithm}

\begin{enumerate}
  \item Global \textit{best} $=-\infty$.
  \item \textsc{Gain}(\textit{node}):
  \begin{enumerate}
    \item If null, return $0$.
    \item $l=\max(0,\textsc{Gain}(\text{left}))$;
          $r=\max(0,\textsc{Gain}(\text{right}))$.
    \item \textit{best} $=\max(\textit{best},\ node.val+l+r)$.
    \item Return $node.val+\max(l,r)$.
  \end{enumerate}
  \item Answer is \textit{best}.
\end{enumerate}

\subsection*{Dry run}

\dryrun{Tree: root $-10$; left $9$; right $20$ with children $15,7$.}

\begin{itemize}
  \item Gain($9$)$=9$, Gain($15$)$=15$, Gain($7$)$=7$.
  \item At $20$: $l=15,r=7$; best$=\max(-\infty,20+15+7)=42$; return
        $20+15=35$.
  \item At $-10$: $l=\max(0,9)=9$, $r=\max(0,35)=35$;
        best$=\max(42,-10+9+35)=42$; return $-10+35=25$.
  \item Maximum path sum $=42$ (path $15\!-\!20\!-\!7$).
\end{itemize}

\subsection*{C++ implementation}

\begin{lstlisting}
#include <bits/stdc++.h>
using namespace std;

struct TreeNode {
    int val;
    TreeNode* left;
    TreeNode* right;
    TreeNode(int v) : val(v), left(nullptr), right(nullptr) {}
};

int best;

int gain(TreeNode* node) {
    if (node == nullptr) return 0;
    int l = max(0, gain(node->left));   // drop harmful subtrees
    int r = max(0, gain(node->right));
    best = max(best, node->val + l + r); // path bending through node
    return node->val + max(l, r);        // parent can continue one way
}

int maxPathSum(TreeNode* root) {
    best = INT_MIN;
    gain(root);
    return best;
}

int main() {
    TreeNode* root = new TreeNode(-10);
    root->left = new TreeNode(9);
    root->right = new TreeNode(20);
    root->right->left = new TreeNode(15);
    root->right->right = new TreeNode(7);

    cout << "Max path sum: " << maxPathSum(root) << "\n";
    return 0;
}
\end{lstlisting}

\begin{complexitybox}
\textbf{Time Complexity.} $O(n)$ -- one postorder pass. \textbf{Space
Complexity.} $O(h)$ for the recursion stack.
\end{complexitybox}

\begin{takeawaybox}
The $\max(0,\cdot)$ clamp -- ``take a subtree's contribution only if it
helps'' -- is the one new idea beyond the diameter template. Whenever a
downward contribution can be negative, clamping it to zero is how greedy
``skip the bad branch'' decisions are expressed inside a bottom-up
recursion.
\end{takeawaybox}

\section{Check if Two Trees Are Identical}
\label{sec:trees-same}

\problemstatement{Given the roots of two binary trees, determine whether
they are structurally identical and have the same node values at every
corresponding position.}

\begin{patternbox}
\emph{``Are these two trees the same?''} is a \textbf{simultaneous
traversal}: walk both trees in lockstep and compare node against node. The
recursion is the same ``ask both children'' shape, but applied to a pair
of nodes at once, with equality being the conjunction of ``roots match''
and ``left subtrees match'' and ``right subtrees match.''
\end{patternbox}

\subsection*{Understanding the problem carefully}

Two trees are identical exactly when their roots hold equal values and,
recursively, their left subtrees are identical and their right subtrees
are identical. The base cases handle the ends of branches: two nulls are
identical; one null and one node are not.

\begin{ideabox}[Lockstep Recursion with a Three-Way Base Case]
Compare node $p$ and node $q$ together:
\begin{itemize}
  \item If both are null: identical here (return true).
  \item If exactly one is null, or their values differ: not identical
        (return false).
  \item Otherwise: return ``left pair identical'' \textbf{and} ``right
        pair identical.''
\end{itemize}
\end{ideabox}

\subsection*{Algorithm}

\begin{enumerate}
  \item \textsc{Same}($p$, $q$):
  \begin{enumerate}
    \item If $p$ and $q$ are both null, return true.
    \item If exactly one is null, return false.
    \item If $p.val \ne q.val$, return false.
    \item Return \textsc{Same}($p$.left, $q$.left) \textbf{and}
          \textsc{Same}($p$.right, $q$.right).
  \end{enumerate}
\end{enumerate}

\subsection*{Dry run}

\dryrun{Tree A: root $1$, left $2$, right $3$. Tree B: root $1$, left $2$,
right $4$.}

\begin{itemize}
  \item Roots $1=1$: recurse on left pair and right pair.
  \item Left pair $2,2$: equal, children all null $\Rightarrow$ true.
  \item Right pair $3,4$: $3\ne4 \Rightarrow$ false.
  \item Conjunction: true \textbf{and} false $=$ \textbf{false}. Not
        identical.
\end{itemize}

\subsection*{C++ implementation}

\begin{lstlisting}
#include <bits/stdc++.h>
using namespace std;

struct TreeNode {
    int val;
    TreeNode* left;
    TreeNode* right;
    TreeNode(int v) : val(v), left(nullptr), right(nullptr) {}
};

bool isSameTree(TreeNode* p, TreeNode* q) {
    if (p == nullptr && q == nullptr) return true;
    if (p == nullptr || q == nullptr) return false;
    if (p->val != q->val) return false;
    return isSameTree(p->left, q->left) &&
           isSameTree(p->right, q->right);
}

int main() {
    TreeNode* a = new TreeNode(1);
    a->left = new TreeNode(2);
    a->right = new TreeNode(3);

    TreeNode* b = new TreeNode(1);
    b->left = new TreeNode(2);
    b->right = new TreeNode(4);

    cout << (isSameTree(a, b) ? "identical" : "different") << "\n";
    return 0;
}
\end{lstlisting}

\begin{complexitybox}
\textbf{Time Complexity.} $O(n)$ where $n$ is the size of the smaller
tree -- the recursion stops at the first mismatch or null. \textbf{Space
Complexity.} $O(h)$ for the recursion stack.
\end{complexitybox}

\begin{takeawaybox}
Comparing two structures by recursing on both in lockstep, with a careful
null/null vs.\ null/non-null base case, is a template that also underlies
``symmetric tree'' (compare a tree against its own mirror) and subtree-of
checks.
\end{takeawaybox}

\section{Zigzag (Spiral) Level Order Traversal}
\label{sec:trees-zigzag}

\problemstatement{Given the root of a binary tree, return its zigzag level
order traversal: the first level left-to-right, the second right-to-left,
the third left-to-right, and so on, alternating direction each level.}

\begin{patternbox}
\emph{``Zigzag,'' ``spiral,'' ``alternate direction each level''} is plain
\textbf{level-order BFS} with one extra bit of state: a flag that flips
every level and controls whether the current level is emitted forwards or
backwards. The traversal structure does not change at all.
\end{patternbox}

\subsection*{Understanding the problem carefully}

BFS already produces nodes grouped by level, left to right. To zigzag, we
only need to reverse the ordering of \emph{every other} level before
adding it to the answer. The cleanest way is to keep the standard
left-to-right BFS but, when the direction flag says ``right-to-left,''
place each node's value into its position from the back of the current
level's list.

\begin{ideabox}[Standard BFS, Reverse Alternate Levels]
Do an ordinary level-order traversal. Maintain a boolean
\textit{leftToRight}. For each level of size $s$, allocate a list of size
$s$; when \textit{leftToRight} is true, write the $i$-th dequeued node at
index $i$; when false, write it at index $s-1-i$. Flip
\textit{leftToRight} after each level.
\end{ideabox}

\subsection*{Algorithm}

\begin{enumerate}
  \item If root is null, return empty. Push root; \textit{leftToRight} $=$
        true.
  \item While the queue is non-empty:
  \begin{enumerate}
    \item $s=$ queue size; make a level array of length $s$.
    \item For $i=0\ldots s-1$: pop node; index $=$ \textit{leftToRight}
          $?\ i : s-1-i$; store value at that index; push children.
    \item Append the level array; flip \textit{leftToRight}.
  \end{enumerate}
\end{enumerate}

\subsection*{Dry run}

\dryrun{Tree: root $1$; children $2,3$; $2$'s children $4,5$; $3$'s
children $6,7$.}

\begin{itemize}
  \item Level $0$ (L$\to$R): $[1]$.
  \item Level $1$ (R$\to$L): nodes dequeued $2,3$ written at indices
        $1,0 \Rightarrow [3,2]$.
  \item Level $2$ (L$\to$R): $[4,5,6,7]$.
  \item Result: $[[1],[3,2],[4,5,6,7]]$.
\end{itemize}

\subsection*{C++ implementation}

\begin{lstlisting}
#include <bits/stdc++.h>
using namespace std;

struct TreeNode {
    int val;
    TreeNode* left;
    TreeNode* right;
    TreeNode(int v) : val(v), left(nullptr), right(nullptr) {}
};

vector<vector<int>> zigzagLevelOrder(TreeNode* root) {
    vector<vector<int>> ans;
    if (root == nullptr) return ans;

    queue<TreeNode*> q;
    q.push(root);
    bool leftToRight = true;
    while (!q.empty()) {
        int s = q.size();
        vector<int> level(s);
        for (int i = 0; i < s; i++) {
            TreeNode* node = q.front();
            q.pop();
            int idx = leftToRight ? i : (s - 1 - i);
            level[idx] = node->val;
            if (node->left)  q.push(node->left);
            if (node->right) q.push(node->right);
        }
        leftToRight = !leftToRight;
        ans.push_back(level);
    }
    return ans;
}

int main() {
    TreeNode* root = new TreeNode(1);
    root->left = new TreeNode(2);
    root->right = new TreeNode(3);
    root->left->left = new TreeNode(4);
    root->left->right = new TreeNode(5);
    root->right->left = new TreeNode(6);
    root->right->right = new TreeNode(7);

    for (auto& lvl : zigzagLevelOrder(root)) {
        for (int x : lvl) cout << x << " ";
        cout << "\n";
    }
    return 0;
}
\end{lstlisting}

\begin{complexitybox}
\textbf{Time Complexity.} $O(n)$ -- one BFS pass; the index arithmetic is
$O(1)$ per node. \textbf{Space Complexity.} $O(n)$ for the queue and
output.
\end{complexitybox}

\begin{takeawaybox}
Many ``fancy-looking'' traversals are ordinary BFS with a single extra
piece of state controlling how each level is read off. Spot the underlying
level-order engine and the decoration -- here, a direction flag -- becomes
trivial.
\end{takeawaybox}

\section{Boundary Traversal of a Binary Tree}
\label{sec:trees-boundary}

\problemstatement{Given the root of a binary tree, return an
anti-clockwise boundary traversal: the left boundary (top-down, excluding
leaves), then all leaves left-to-right, then the right boundary
(bottom-up, excluding leaves). The root is included once.}

\begin{patternbox}
\emph{``Boundary,'' ``outline,'' ``perimeter''} decomposes into three
independent walks that must be stitched without double-counting: the left
edge going down, the leaves going across, and the right edge coming back
up. The only real care is avoiding overlap where these parts meet -- which
is why leaves are deliberately excluded from the two edges.
\end{patternbox}

\subsection*{Understanding the problem carefully}

The perimeter of the tree consists of three disjoint pieces. The
\textbf{left boundary} is the path you get by always going left (or right
if there is no left child), from the root's left child down, stopping
before leaves. The \textbf{leaves} are every node with no children, in
left-to-right order. The \textbf{right boundary} is the mirror -- always
go right (or left otherwise) -- collected then reversed so it reads
bottom-up. Excluding leaves from the two boundaries prevents a corner leaf
from being listed twice.

\begin{ideabox}[Three Walks, No Overlap]
(1) Add the root (unless it is itself a leaf). (2) Walk the left boundary
top-down, adding non-leaf nodes. (3) Add all leaves via any DFS,
left-to-right. (4) Walk the right boundary top-down into a temporary list,
adding non-leaf nodes, then append that list \emph{reversed}. Leaves are
skipped in steps~2 and~4 precisely so step~3 owns them exclusively.
\end{ideabox}

\subsection*{Algorithm}

\begin{enumerate}
  \item If root is not a leaf, add root's value.
  \item Left boundary: start at root.left; while the node exists, if it is
        not a leaf add it; move to its left child, or right child if no
        left.
  \item Leaves: DFS the whole tree; whenever a node has no children, add
        it.
  \item Right boundary: start at root.right; collect non-leaf nodes
        top-down (preferring right children) into a list; append that
        list reversed.
\end{enumerate}

\subsection*{Dry run}

\dryrun{Tree: root $1$; left $2$ (children $4$,$5$); right $3$ (right
child $6$, whose left child $7$).}

\begin{itemize}
  \item Root $1$ (not a leaf): add $1$.
  \item Left boundary from $2$: $2$ not a leaf, add $2$; go to $4$ -- a
        leaf, stop. Boundary so far: $1,2$.
  \item Leaves L$\to$R: $4,5,7$. Now: $1,2,4,5,7$.
  \item Right boundary from $3$: $3$ not leaf add $3$; go to $6$, not leaf
        add $6$; $7$ is a leaf, stop. List $[3,6]$ reversed $=[6,3]$.
  \item Result: $1,2,4,5,7,6,3$.
\end{itemize}

\subsection*{C++ implementation}

\begin{lstlisting}
#include <bits/stdc++.h>
using namespace std;

struct TreeNode {
    int val;
    TreeNode* left;
    TreeNode* right;
    TreeNode(int v) : val(v), left(nullptr), right(nullptr) {}
};

bool isLeaf(TreeNode* n) {
    return n->left == nullptr && n->right == nullptr;
}

void addLeaves(TreeNode* n, vector<int>& out) {
    if (n == nullptr) return;
    if (isLeaf(n)) { out.push_back(n->val); return; }
    addLeaves(n->left, out);
    addLeaves(n->right, out);
}

vector<int> boundaryTraversal(TreeNode* root) {
    vector<int> out;
    if (root == nullptr) return out;
    if (!isLeaf(root)) out.push_back(root->val);

    // left boundary (top-down, non-leaf)
    TreeNode* cur = root->left;
    while (cur) {
        if (!isLeaf(cur)) out.push_back(cur->val);
        cur = cur->left ? cur->left : cur->right;
    }

    // all leaves, left to right
    addLeaves(root, out);

    // right boundary (top-down into temp, then reversed)
    vector<int> tmp;
    cur = root->right;
    while (cur) {
        if (!isLeaf(cur)) tmp.push_back(cur->val);
        cur = cur->right ? cur->right : cur->left;
    }
    for (int i = (int)tmp.size() - 1; i >= 0; i--) out.push_back(tmp[i]);
    return out;
}

int main() {
    TreeNode* root = new TreeNode(1);
    root->left = new TreeNode(2);
    root->left->left = new TreeNode(4);
    root->left->right = new TreeNode(5);
    root->right = new TreeNode(3);
    root->right->right = new TreeNode(6);
    root->right->right->left = new TreeNode(7);

    for (int x : boundaryTraversal(root)) cout << x << " ";
    cout << "\n";
    return 0;
}
\end{lstlisting}

\begin{complexitybox}
\textbf{Time Complexity.} $O(n)$ -- the leaf DFS is $O(n)$ and each
boundary walk is $O(h)$. \textbf{Space Complexity.} $O(n)$ for the output
and $O(h)$ recursion for the leaf pass.
\end{complexitybox}

\begin{takeawaybox}
Boundary problems are less about clever algorithms and more about a clean
decomposition with disciplined de-duplication at the seams. Deciding up
front that ``leaves belong to exactly one of the three parts'' is what
keeps the corners from being counted twice.
\end{takeawaybox}

\section{Vertical Order Traversal}
\label{sec:trees-vertical}

\problemstatement{Assign the root a horizontal distance (HD) of $0$; a left
child has HD one less than its parent, a right child one more. Return
nodes grouped by HD from leftmost to rightmost. Within the same HD, order
by depth (top to bottom); nodes at the same HD \emph{and} same depth are
ordered by value.}

\begin{patternbox}
\emph{``Vertical,'' ``columns,'' ``horizontal distance''} says: tag each
node with a \textbf{(HD, depth)} coordinate, collect everything into a map
keyed by HD, and sort within each column. The traversal is ordinary BFS or
DFS; the work is in the coordinate bookkeeping and the tie-breaking rules.
\end{patternbox}

\subsection*{Understanding the problem carefully}

Imagine dropping vertical lines through the tree. Every node falls on one
line, identified by its horizontal distance. Reading the answer means
listing lines left to right, and within a line, top to bottom. The precise
LeetCode-style rule adds a final tie-break: if two nodes share both HD and
depth, the smaller value comes first. Storing each node as a triple
$(\text{HD},\text{depth},\text{value})$ and sorting makes all three rules
fall out of one comparison.

\begin{ideabox}[Order by (HD, then depth, then value)]
Traverse the tree recording, for each node, the triple $(HD, depth,
val)$: root is $(0,0,val)$, a left child inherits $(HD-1, depth+1)$, a
right child $(HD+1, depth+1)$. Group into a map from HD to the list of
$(depth, val)$ pairs; sort each list by $(depth, val)$; emit columns in
increasing HD.
\end{ideabox}

\subsection*{Algorithm}

\begin{enumerate}
  \item BFS/DFS from the root carrying $(HD, depth)$; append $(depth,
        val)$ into $\textit{cols}[HD]$.
  \item For each HD in increasing order, sort its list by $(depth, val)$
        and output the values.
\end{enumerate}

\subsection*{Dry run}

\dryrun{Tree: root $1$; left $2$; right $3$; $2$'s right child $4$; $3$'s
left child $5$ (so $4$ and $5$ both land on HD $0$).}

\begin{itemize}
  \item $1:(0,0)$; $2:(-1,1)$; $3:(1,1)$; $4:(0,2)$; $5:(0,2)$.
  \item HD $-1$: $[2]$. HD $0$: $\{(0,1),(2,4),(2,5)\}$ sorted by
        (depth,val) $\to 1,4,5$. HD $1$: $[3]$.
  \item Result: $[[2],[1,4,5],[3]]$.
\end{itemize}

\subsection*{C++ implementation}

\begin{lstlisting}
#include <bits/stdc++.h>
using namespace std;

struct TreeNode {
    int val;
    TreeNode* left;
    TreeNode* right;
    TreeNode(int v) : val(v), left(nullptr), right(nullptr) {}
};

void dfs(TreeNode* node, int hd, int depth,
         map<int, vector<pair<int,int>>>& cols) {
    if (node == nullptr) return;
    cols[hd].push_back({depth, node->val});
    dfs(node->left,  hd - 1, depth + 1, cols);
    dfs(node->right, hd + 1, depth + 1, cols);
}

vector<vector<int>> verticalTraversal(TreeNode* root) {
    map<int, vector<pair<int,int>>> cols;   // hd -> (depth, val)
    dfs(root, 0, 0, cols);

    vector<vector<int>> ans;
    for (auto& [hd, vec] : cols) {          // map iterates hd ascending
        sort(vec.begin(), vec.end());       // by depth, then value
        vector<int> col;
        for (auto& [d, v] : vec) col.push_back(v);
        ans.push_back(col);
    }
    return ans;
}

int main() {
    TreeNode* root = new TreeNode(1);
    root->left = new TreeNode(2);
    root->right = new TreeNode(3);
    root->left->right = new TreeNode(4);
    root->right->left = new TreeNode(5);

    for (auto& col : verticalTraversal(root)) {
        for (int x : col) cout << x << " ";
        cout << "\n";
    }
    return 0;
}
\end{lstlisting}

\begin{complexitybox}
\textbf{Time Complexity.} $O(n\log n)$ -- dominated by sorting the columns
(total $n$ entries across all columns). \textbf{Space Complexity.} $O(n)$
for the map and the recursion stack.
\end{complexitybox}

\begin{takeawaybox}
Tagging nodes with coordinates and offloading the ordering to a sorted map
plus a per-column sort turns a geometric-sounding problem into
bookkeeping. Top view and bottom view (next sections) reuse the same HD
tag with a simpler rule for which node on each line ``wins.''
\end{takeawaybox}

\section{Top View of a Binary Tree}
\label{sec:trees-top-view}

\problemstatement{The top view is the set of nodes visible when the tree is
viewed from directly above: for each horizontal distance (HD), the single
\emph{topmost} node on that vertical line. Return these nodes ordered left
to right by HD.}

\begin{patternbox}
Reusing the HD tag from vertical order, \emph{``top view / viewed from
above''} means: on each vertical line, keep only the node encountered
\emph{first when scanning top-down}. A \textbf{BFS} scans strictly in
increasing depth, so the first node BFS places on a given HD is
automatically the topmost -- record an HD only the first time it appears.
\end{patternbox}

\subsection*{Understanding the problem carefully}

Looking down from above, a node is hidden if some other node shares its
horizontal distance but sits higher. Because BFS processes level $0$
before level $1$ before level $2$, the earliest time an HD is seen during
BFS corresponds to the highest node on that line. So we simply record the
first value per HD and ignore all later ones.

\begin{ideabox}[BFS, First Node per HD Wins]
Run level-order BFS carrying each node's HD (root $0$, left $-1$, right
$+1$). Keep a map from HD to value, writing an entry \emph{only if that HD
is not already present}. BFS's increasing-depth order guarantees the first
write for each HD is the topmost node. Emit the map in ascending HD order.
\end{ideabox}

\emph{Why BFS, not DFS:} a DFS might reach a deep node on some HD before a
shallower node on the same HD (down a different branch), wrongly recording
the deeper one. BFS's level-by-level order removes that hazard.

\subsection*{Algorithm}

\begin{enumerate}
  \item BFS from root with a queue of $(node, HD)$; root's HD $=0$.
  \item For each dequeued $(node, HD)$: if $HD$ not yet in the map, set
        $\textit{map}[HD]=node.val$. Enqueue $(left, HD-1)$ and
        $(right, HD+1)$.
  \item Output the map's values in ascending HD.
\end{enumerate}

\subsection*{Dry run}

\dryrun{Tree: root $1$; left $2$ (right child $4$); right $3$.}

\begin{itemize}
  \item $(1,0)$: map$\{0\!:\!1\}$; enqueue $(2,-1),(3,1)$.
  \item $(2,-1)$: map$\{-1\!:\!2,0\!:\!1\}$; enqueue $(4,0)$.
  \item $(3,1)$: map adds $1\!:\!3$.
  \item $(4,0)$: HD $0$ already present -- skip.
  \item Ascending HD: $2,1,3$.
\end{itemize}

\subsection*{C++ implementation}

\begin{lstlisting}
#include <bits/stdc++.h>
using namespace std;

struct TreeNode {
    int val;
    TreeNode* left;
    TreeNode* right;
    TreeNode(int v) : val(v), left(nullptr), right(nullptr) {}
};

vector<int> topView(TreeNode* root) {
    vector<int> ans;
    if (root == nullptr) return ans;

    map<int, int> topAtHD;               // hd -> first (topmost) value
    queue<pair<TreeNode*, int>> q;
    q.push({root, 0});
    while (!q.empty()) {
        auto [node, hd] = q.front();
        q.pop();
        if (topAtHD.find(hd) == topAtHD.end())
            topAtHD[hd] = node->val;     // first sighting wins
        if (node->left)  q.push({node->left,  hd - 1});
        if (node->right) q.push({node->right, hd + 1});
    }
    for (auto& [hd, v] : topAtHD) ans.push_back(v);
    return ans;
}

int main() {
    TreeNode* root = new TreeNode(1);
    root->left = new TreeNode(2);
    root->right = new TreeNode(3);
    root->left->right = new TreeNode(4);

    for (int x : topView(root)) cout << x << " ";
    cout << "\n";
    return 0;
}
\end{lstlisting}

\begin{complexitybox}
\textbf{Time Complexity.} $O(n\log n)$ -- $n$ BFS steps, each doing an
$O(\log n)$ map operation (or $O(n)$ overall with a hash map plus a final
sort). \textbf{Space Complexity.} $O(n)$ for the queue and map.
\end{complexitybox}

\begin{takeawaybox}
``First per HD under BFS'' is the whole trick for top view. The insistence
on BFS over DFS is the subtle correctness point: only a strictly
increasing-depth scan makes ``first seen'' equal to ``topmost.''
\end{takeawaybox}

\section{Bottom View of a Binary Tree}
\label{sec:trees-bottom-view}

\problemstatement{The bottom view is the set of nodes visible when the tree
is viewed from directly below: for each horizontal distance (HD), the
\emph{lowest} node on that vertical line (ties broken by the later node in
level order). Return these ordered left to right by HD.}

\begin{patternbox}
Bottom view is top view with one word changed: keep the \emph{last} node
per HD instead of the first. Under BFS -- which scans top-down -- the last
node written to a given HD is the deepest, exactly the one visible from
below. So the code is identical except we \textbf{always overwrite} the
map entry rather than writing only once.
\end{patternbox}

\subsection*{Understanding the problem carefully}

From below, on each vertical line the visible node is the lowest one.
Because BFS visits nodes in non-decreasing depth, repeatedly overwriting
the value stored for an HD leaves, at the end, the value of the last (hence
deepest) node BFS placed there. The left-to-right order within a level
also gives the standard tie-break for two nodes at the same HD and depth.

\begin{ideabox}[BFS, Last Node per HD Wins]
Run the same level-order BFS with HD tags as top view, but write
$\textit{map}[HD]=node.val$ \emph{unconditionally} every time. Later
(deeper) nodes overwrite earlier ones, so each HD ends holding its
bottommost node. Emit ascending by HD.
\end{ideabox}

\subsection*{Algorithm}

\begin{enumerate}
  \item BFS from root with $(node, HD)$; root's HD $=0$.
  \item For each dequeued $(node, HD)$: set $\textit{map}[HD]=node.val$
        (overwrite). Enqueue $(left, HD-1)$, $(right, HD+1)$.
  \item Output the map's values in ascending HD.
\end{enumerate}

\subsection*{Dry run}

\dryrun{Tree: root $1$; left $2$ (right child $4$); right $3$.}

\begin{itemize}
  \item $(1,0)$: map$\{0\!:\!1\}$; enqueue $(2,-1),(3,1)$.
  \item $(2,-1)$: map$\{-1\!:\!2\}$; enqueue $(4,0)$.
  \item $(3,1)$: map adds $1\!:\!3$.
  \item $(4,0)$: overwrite HD $0$ with $4$.
  \item Ascending HD: $2,4,3$ (note HD $0$ now shows $4$, not $1$).
\end{itemize}

\subsection*{C++ implementation}

\begin{lstlisting}
#include <bits/stdc++.h>
using namespace std;

struct TreeNode {
    int val;
    TreeNode* left;
    TreeNode* right;
    TreeNode(int v) : val(v), left(nullptr), right(nullptr) {}
};

vector<int> bottomView(TreeNode* root) {
    vector<int> ans;
    if (root == nullptr) return ans;

    map<int, int> lowAtHD;               // hd -> last (lowest) value
    queue<pair<TreeNode*, int>> q;
    q.push({root, 0});
    while (!q.empty()) {
        auto [node, hd] = q.front();
        q.pop();
        lowAtHD[hd] = node->val;         // always overwrite
        if (node->left)  q.push({node->left,  hd - 1});
        if (node->right) q.push({node->right, hd + 1});
    }
    for (auto& [hd, v] : lowAtHD) ans.push_back(v);
    return ans;
}

int main() {
    TreeNode* root = new TreeNode(1);
    root->left = new TreeNode(2);
    root->right = new TreeNode(3);
    root->left->right = new TreeNode(4);

    for (int x : bottomView(root)) cout << x << " ";
    cout << "\n";
    return 0;
}
\end{lstlisting}

\begin{complexitybox}
\textbf{Time Complexity.} $O(n\log n)$ with an ordered map (or $O(n)$ with
a hash map and a final sort). \textbf{Space Complexity.} $O(n)$ for the
queue and map.
\end{complexitybox}

\begin{takeawaybox}
Top view and bottom view differ by exactly one decision -- keep the first
versus the last node per HD. Recognising that a whole family of ``view''
problems shares one BFS-with-HD engine, differing only in the per-line
selection rule, is the real lesson.
\end{takeawaybox}

\section{Left View and Right View of a Binary Tree}
\label{sec:trees-left-right-view}

\problemstatement{The right view is the sequence of nodes visible from the
right side of the tree: the \emph{last} node at each level (top to bottom).
The left view is the \emph{first} node at each level. Return both.}

\begin{patternbox}
\emph{``View from the side,'' ``first/last node per level''} is level-order
selection: for each level, the rightmost node is the right view, the
leftmost is the left view. This is even simpler than the HD-based views --
there is no horizontal distance, only ``which node ends (or begins) each
level.''
\end{patternbox}

\subsection*{Understanding the problem carefully}

Standing to the right of the tree, on each level you see only the node
furthest right, because it blocks the ones behind it. So the right view is
just ``the last node of each level.'' Equivalently, with a depth-first
approach that visits \emph{right child before left}, the first node seen
at each new depth is the right-view node -- which gives an elegant $O(h)$
recursion.

\begin{ideabox}[First Node Seen per Depth, Right-First DFS]
Recurse carrying the current depth. Visit the \emph{right} child before
the left. The first time the recursion reaches a given depth, the node it
is at is the rightmost at that depth -- record it. For the left view, do
the same but visit the \emph{left} child first.
\end{ideabox}

\subsection*{Algorithm (right view)}

\begin{enumerate}
  \item \textsc{DFS}(\textit{node}, \textit{depth}, \textit{out}):
  \begin{enumerate}
    \item If \textit{node} is null, return.
    \item If $\textit{depth} = \text{size of } \textit{out}$: this is the
          first node reached at this depth -- append its value.
    \item \textsc{DFS}(\textit{node}.right, \textit{depth}+1);
          \textsc{DFS}(\textit{node}.left, \textit{depth}+1).
  \end{enumerate}
\end{enumerate}

\subsection*{Dry run}

\dryrun{Tree: root $1$; left $2$ (right child $5$); right $3$.}

\begin{itemize}
  \item DFS($1$,$0$): out empty, depth $0=0 \Rightarrow$ record $1$.
        Recurse right into $3$, then left into $2$.
  \item DFS($3$,$1$): depth $1=$ size $1 \Rightarrow$ record $3$. (No
        children.)
  \item DFS($2$,$1$): depth $1 \ne$ size $2$ -- skip. Recurse right into
        $5$.
  \item DFS($5$,$2$): depth $2=$ size $2 \Rightarrow$ record $5$.
  \item Right view: $1,3,5$. (Left view by symmetry: $1,2,5$.)
\end{itemize}

\subsection*{C++ implementation}

\begin{lstlisting}
#include <bits/stdc++.h>
using namespace std;

struct TreeNode {
    int val;
    TreeNode* left;
    TreeNode* right;
    TreeNode(int v) : val(v), left(nullptr), right(nullptr) {}
};

void rightDFS(TreeNode* node, int depth, vector<int>& out) {
    if (node == nullptr) return;
    if (depth == (int)out.size()) out.push_back(node->val);
    rightDFS(node->right, depth + 1, out);   // right first
    rightDFS(node->left,  depth + 1, out);
}

void leftDFS(TreeNode* node, int depth, vector<int>& out) {
    if (node == nullptr) return;
    if (depth == (int)out.size()) out.push_back(node->val);
    leftDFS(node->left,  depth + 1, out);    // left first
    leftDFS(node->right, depth + 1, out);
}

int main() {
    TreeNode* root = new TreeNode(1);
    root->left = new TreeNode(2);
    root->right = new TreeNode(3);
    root->left->right = new TreeNode(5);

    vector<int> rv, lv;
    rightDFS(root, 0, rv);
    leftDFS(root, 0, lv);

    cout << "Right view: "; for (int x : rv) cout << x << " "; cout << "\n";
    cout << "Left view:  "; for (int x : lv) cout << x << " "; cout << "\n";
    return 0;
}
\end{lstlisting}

\begin{complexitybox}
\textbf{Time Complexity.} $O(n)$ -- every node is visited once.
\textbf{Space Complexity.} $O(h)$ for the recursion stack.
\end{complexitybox}

\begin{takeawaybox}
The ``record only when depth equals the current answer size'' test is a
compact way to detect ``first node reached at a new depth'' without an
explicit level loop. Swapping the child-visit order flips left view to
right view -- another example of one algorithm, two mirrored variants.
\end{takeawaybox}

\section{Symmetric Binary Tree}
\label{sec:trees-symmetric}

\problemstatement{A binary tree is symmetric if it is a mirror image of
itself around its centre: the left subtree is the mirror reflection of the
right subtree. Given the root, determine whether the tree is symmetric.}

\begin{patternbox}
\emph{``Mirror,'' ``symmetric,'' ``reflection''} is the identical-trees
lockstep recursion (Section~\ref{sec:trees-same}) with a twist: instead of
comparing left-with-left and right-with-right, we compare the \emph{outer}
pair with each other and the \emph{inner} pair with each other. Mirror
matching crosses the two children.
\end{patternbox}

\subsection*{Understanding the problem carefully}

Symmetry is a statement about the left subtree versus the right subtree.
Two subtrees are mirror images when their roots are equal \emph{and} the
left one's left child mirrors the right one's right child (the outer edges)
\emph{and} the left one's right child mirrors the right one's left child
(the inner edges). The crossing -- left.left with right.right, left.right
with right.left -- is what encodes ``reflection'' rather than ``equality.''

\begin{ideabox}[Mirror-Compare: Outer with Outer, Inner with Inner]
Define \textsc{Mirror}($a$, $b$): they mirror iff both are null, or both
are non-null with equal values \emph{and} \textsc{Mirror}($a$.left,
$b$.right) \emph{and} \textsc{Mirror}($a$.right, $b$.left). The whole tree
is symmetric iff \textsc{Mirror}(root.left, root.right).
\end{ideabox}

\subsection*{Algorithm}

\begin{enumerate}
  \item If root is null, return true.
  \item \textsc{Mirror}($a$, $b$):
  \begin{enumerate}
    \item If both null, return true; if exactly one null, return false.
    \item If $a.val \ne b.val$, return false.
    \item Return \textsc{Mirror}($a$.left, $b$.right) \textbf{and}
          \textsc{Mirror}($a$.right, $b$.left).
  \end{enumerate}
  \item Return \textsc{Mirror}(root.left, root.right).
\end{enumerate}

\subsection*{Dry run}

\dryrun{Tree: root $1$; left $2$ (children $3,4$); right $2$ (children
$4,3$).}

\begin{itemize}
  \item Mirror($2_L$, $2_R$): values equal. Compare outer
        ($2_L$.left$=3$, $2_R$.right$=3$) and inner ($2_L$.right$=4$,
        $2_R$.left$=4$).
  \item Outer $3,3$: equal, children null $\to$ true. Inner $4,4$: true.
  \item All true $\Rightarrow$ \textbf{symmetric}.
\end{itemize}

\subsection*{C++ implementation}

\begin{lstlisting}
#include <bits/stdc++.h>
using namespace std;

struct TreeNode {
    int val;
    TreeNode* left;
    TreeNode* right;
    TreeNode(int v) : val(v), left(nullptr), right(nullptr) {}
};

bool mirror(TreeNode* a, TreeNode* b) {
    if (a == nullptr && b == nullptr) return true;
    if (a == nullptr || b == nullptr) return false;
    if (a->val != b->val) return false;
    return mirror(a->left, b->right) &&    // outer pair
           mirror(a->right, b->left);      // inner pair
}

bool isSymmetric(TreeNode* root) {
    if (root == nullptr) return true;
    return mirror(root->left, root->right);
}

int main() {
    TreeNode* root = new TreeNode(1);
    root->left = new TreeNode(2);
    root->right = new TreeNode(2);
    root->left->left = new TreeNode(3);
    root->left->right = new TreeNode(4);
    root->right->left = new TreeNode(4);
    root->right->right = new TreeNode(3);

    cout << (isSymmetric(root) ? "symmetric" : "not symmetric") << "\n";
    return 0;
}
\end{lstlisting}

\begin{complexitybox}
\textbf{Time Complexity.} $O(n)$ -- each pair of mirrored nodes is compared
once. \textbf{Space Complexity.} $O(h)$ for the recursion stack.
\end{complexitybox}

\begin{takeawaybox}
Symmetry is equality with the two children crossed. Recognising it as ``the
same-tree recursion, but mirror the child pairing'' means you write it
correctly the first time instead of tangling the four child comparisons.
\end{takeawaybox}


% ---- Hard ----
\section{Root-to-Node Path}
\label{sec:trees-root-to-node-path}

\problemstatement{Given the root of a binary tree and a target node value
(assume values are unique), return the sequence of node values on the path
from the root down to the target.}

\begin{patternbox}
\emph{``Path from the root to a node,'' ``ancestors of a node''} is a
backtracking DFS: descend adding the current node to a running path; if
this node is the target, stop and keep the path; if neither subtree
contains the target, \textbf{remove the current node} before returning. The
add-then-remove-on-failure rhythm is the essence of backtracking.
\end{patternbox}

\subsection*{Understanding the problem carefully}

There is exactly one path from the root to any node in a tree. We build it
incrementally: push the current node onto the path, then check ``is this
the target, or is the target somewhere below on the left or right?'' If
yes, this node is genuinely on the path and stays. If the target is not
here and not in either subtree, this node was a wrong turn and must be
popped so it does not pollute the path for other branches.

\begin{ideabox}[Add on Entry, Pop on a Dead End]
\textsc{Find}(node, target, path): if node is null return false; push
node.val onto path; if node.val equals target return true; if
\textsc{Find}(left) or \textsc{Find}(right) returns true, return true
(node stays on the path); otherwise pop node.val (dead end) and return
false.
\end{ideabox}

\subsection*{Algorithm}

\begin{enumerate}
  \item \textsc{Find}(\textit{node}, \textit{target}, \textit{path}):
  \begin{enumerate}
    \item If \textit{node} is null, return false.
    \item Push \textit{node}.val onto \textit{path}.
    \item If \textit{node}.val $=$ \textit{target}, return true.
    \item If \textsc{Find}(left) \textbf{or} \textsc{Find}(right),
          return true.
    \item Pop \textit{node}.val; return false.
  \end{enumerate}
\end{enumerate}

\subsection*{Dry run}

\dryrun{Tree: root $1$; left $2$ (children $4,5$); right $3$. Target $5$.}

\begin{itemize}
  \item Push $1$. Not $5$. Recurse left into $2$.
  \item Push $2$. Not $5$. Recurse left into $4$: push $4$, not $5$, no
        children, pop $4$, return false. Recurse right into $5$: push
        $5$, equals target, return true.
  \item $2$ keeps ($5$ found below); $1$ keeps.
  \item Path: $1,2,5$.
\end{itemize}

\subsection*{C++ implementation}

\begin{lstlisting}
#include <bits/stdc++.h>
using namespace std;

struct TreeNode {
    int val;
    TreeNode* left;
    TreeNode* right;
    TreeNode(int v) : val(v), left(nullptr), right(nullptr) {}
};

bool findPath(TreeNode* node, int target, vector<int>& path) {
    if (node == nullptr) return false;
    path.push_back(node->val);
    if (node->val == target) return true;
    if (findPath(node->left, target, path) ||
        findPath(node->right, target, path))
        return true;
    path.pop_back();                 // dead end: backtrack
    return false;
}

int main() {
    TreeNode* root = new TreeNode(1);
    root->left = new TreeNode(2);
    root->right = new TreeNode(3);
    root->left->left = new TreeNode(4);
    root->left->right = new TreeNode(5);

    vector<int> path;
    findPath(root, 5, path);
    for (int x : path) cout << x << " ";
    cout << "\n";
    return 0;
}
\end{lstlisting}

\begin{complexitybox}
\textbf{Time Complexity.} $O(n)$ -- in the worst case every node is
visited once. \textbf{Space Complexity.} $O(h)$ for the recursion stack
and the path.
\end{complexitybox}

\begin{takeawaybox}
The push-on-entry / pop-on-failure pattern is backtracking in its simplest
tree form, and it is the conceptual basis for the next section's lowest
common ancestor and for any ``list the ancestors'' query.
\end{takeawaybox}

\section{Lowest Common Ancestor in a Binary Tree}
\label{sec:trees-lca}

\problemstatement{Given the root of a binary tree and two nodes $p$ and
$q$ (both guaranteed present), return their lowest common ancestor: the
deepest node that has both $p$ and $q$ as descendants (a node is a
descendant of itself).}

\begin{patternbox}
\emph{``Lowest common ancestor'' in a \textbf{general} binary tree} (no BST
ordering to exploit) is a single bottom-up recursion that reports back up
which of $p$ and $q$ each subtree contains. The LCA is the unique node
where the two targets are found on \emph{different sides} -- or which is
itself one target with the other below it.
\end{patternbox}

\subsection*{Understanding the problem carefully}

Walk the tree once. Each recursive call answers: ``does my subtree contain
$p$ or $q$ (or one of them being me)?'' by returning a non-null pointer if
so. If a node finds one target in its left subtree and the other in its
right subtree, that node is the meeting point -- the LCA. If both are found
on the same side, the LCA lies deeper on that side and is propagated up
unchanged.

\begin{ideabox}[The Split Point Is the Answer]
\textsc{LCA}(node): if node is null or equals $p$ or $q$, return node.
Recurse left and right. If \emph{both} calls return non-null, the current
node is where the two targets split -- return it. Otherwise return whichever
side was non-null (or null). The topmost node receiving non-null from both
sides is the LCA.
\end{ideabox}

\subsection*{Algorithm}

\begin{enumerate}
  \item \textsc{LCA}(\textit{node}):
  \begin{enumerate}
    \item If \textit{node} is null or equals $p$ or $q$, return
          \textit{node}.
    \item $L=$\textsc{LCA}(left); $R=$\textsc{LCA}(right).
    \item If $L$ and $R$ are both non-null, return \textit{node}.
    \item Return $L$ if non-null, else $R$.
  \end{enumerate}
\end{enumerate}

\subsection*{Dry run}

\dryrun{Tree: root $3$; left $5$ (children $6,2$); right $1$. Find LCA of
$6$ and $2$.}

\begin{itemize}
  \item At $5$: left call at $6$ returns $6$; right call at $2$ returns
        $2$. Both non-null $\Rightarrow$ $5$ is the LCA; return $5$.
  \item At $3$: left returns $5$, right (subtree $1$) returns null. Only
        left non-null $\Rightarrow$ return $5$.
  \item LCA $=5$.
\end{itemize}

\subsection*{C++ implementation}

\begin{lstlisting}
#include <bits/stdc++.h>
using namespace std;

struct TreeNode {
    int val;
    TreeNode* left;
    TreeNode* right;
    TreeNode(int v) : val(v), left(nullptr), right(nullptr) {}
};

TreeNode* lca(TreeNode* node, TreeNode* p, TreeNode* q) {
    if (node == nullptr || node == p || node == q) return node;
    TreeNode* L = lca(node->left,  p, q);
    TreeNode* R = lca(node->right, p, q);
    if (L != nullptr && R != nullptr) return node;  // split point
    return L != nullptr ? L : R;
}

int main() {
    TreeNode* root = new TreeNode(3);
    root->left = new TreeNode(5);
    root->right = new TreeNode(1);
    TreeNode* p = root->left->left  = new TreeNode(6);
    TreeNode* q = root->left->right = new TreeNode(2);

    cout << "LCA: " << lca(root, p, q)->val << "\n";
    return 0;
}
\end{lstlisting}

\begin{complexitybox}
\textbf{Time Complexity.} $O(n)$ -- a single traversal. \textbf{Space
Complexity.} $O(h)$ for the recursion stack.
\end{complexitybox}

\begin{takeawaybox}
Returning a ``found'' pointer up the recursion and declaring the LCA at
the node where left and right both report success is a one-pass gem. Unlike
the BST version, no value ordering is used -- only the tree's shape --
which is why it works on any binary tree.
\end{takeawaybox}

\section{Maximum Width of a Binary Tree}
\label{sec:trees-max-width}

\problemstatement{The width of a level is the number of positions between
its leftmost and rightmost non-null nodes, \emph{including} the null gaps
between them (as if the tree were a complete binary tree). Return the
maximum width over all levels.}

\begin{patternbox}
The phrase \emph{``including null gaps''} is the whole difficulty:
counting only present nodes is wrong. The fix is to give every node a
\textbf{positional index} as if the tree were complete -- a node at index
$i$ has children at $2i$ and $2i+1$ -- and measure each level's width as
(rightmost index $-$ leftmost index $+1$). This is BFS with indices.
\end{patternbox}

\subsection*{Understanding the problem carefully}

In a complete binary tree numbered from the root as index $0$, a node at
index $i$ places its left child at $2i+1$ and right child at $2i+2$. These
indices encode horizontal position \emph{including} the empty slots. So for
each level, the first and last real nodes carry indices whose difference,
plus one, is exactly the width that counts the gaps between them. To keep
the indices from overflowing on deep trees, subtract the level's leftmost
index from every node's index (re-basing each level to start at $0$).

\begin{ideabox}[Number Nodes as if the Tree Were Complete]
BFS level by level, pairing each node with an index (root $=0$; children
$2i+1$, $2i+2$). For each level, width $=$ (last index $-$ first index
$+1$); track the maximum. Re-base each level's indices by subtracting the
first node's index to prevent overflow.
\end{ideabox}

\subsection*{Algorithm}

\begin{enumerate}
  \item Queue holds $(node, index)$; push $(root, 0)$.
  \item For each level: let \textit{first} be the index of the level's
        first node. For every node, push $(left, 2\cdot(i-\textit{first})+1)$
        and $(right, 2\cdot(i-\textit{first})+2)$. Width $=$ last index
        $-$ first index $+1$; update the maximum.
\end{enumerate}

\subsection*{Dry run}

\dryrun{Tree: root $1$; left $2$ (left child $4$); right $3$ (right child
$5$).}

\begin{itemize}
  \item Level $0$: $(1,0)$. Width $1$.
  \item Level $1$: $(2,0),(3,1)$ (re-based). Width $1-0+1=2$.
  \item Level $2$: from $2$ push $(4,1)$; from $3$ push $(5,4)$. First
        index $1$; re-base to $(4,0),(5,3)$. Width $3-0+1=4$.
  \item Maximum width $=4$.
\end{itemize}

\subsection*{C++ implementation}

\begin{lstlisting}
#include <bits/stdc++.h>
using namespace std;

struct TreeNode {
    int val;
    TreeNode* left;
    TreeNode* right;
    TreeNode(int v) : val(v), left(nullptr), right(nullptr) {}
};

int widthOfBinaryTree(TreeNode* root) {
    if (root == nullptr) return 0;
    long long best = 0;
    queue<pair<TreeNode*, long long>> q;
    q.push({root, 0});
    while (!q.empty()) {
        int s = q.size();
        long long first = q.front().second, last = first;
        for (int i = 0; i < s; i++) {
            auto [node, idx] = q.front();
            q.pop();
            idx -= first;                 // re-base to avoid overflow
            last = idx;
            if (node->left)  q.push({node->left,  2 * idx + 1});
            if (node->right) q.push({node->right, 2 * idx + 2});
        }
        best = max(best, last + 1);
    }
    return (int)best;
}

int main() {
    TreeNode* root = new TreeNode(1);
    root->left = new TreeNode(2);
    root->right = new TreeNode(3);
    root->left->left = new TreeNode(4);
    root->right->right = new TreeNode(5);

    cout << "Max width: " << widthOfBinaryTree(root) << "\n";
    return 0;
}
\end{lstlisting}

\begin{complexitybox}
\textbf{Time Complexity.} $O(n)$ -- one BFS pass. \textbf{Space
Complexity.} $O(n)$ for the queue.
\end{complexitybox}

\begin{takeawaybox}
Positional indexing ($2i+1$, $2i+2$) is the standard way to reason about a
tree's geometry when empty slots matter. Re-basing per level is the
practical detail that keeps those indices from overflowing on a deep,
sparse tree.
\end{takeawaybox}

\section{Children Sum Property}
\label{sec:trees-children-sum}

\problemstatement{The children-sum property holds if, for every node with
at least one child, the node's value equals the sum of its children's
values (a missing child counts as $0$). Given a tree, \emph{modify node
values} (you may only increase them) so the property holds, without
changing the tree's structure.}

\begin{patternbox}
This is not a check but a \textbf{construction}, and the ``you may only
increase values'' rule forces a specific two-pass shape: push required
values \emph{down} on the way in (raise children up to at least the
parent), then fix parents \emph{up} on the way out (set each parent to its
children's actual sum). The down-then-up structure is the key insight.
\end{patternbox}

\subsection*{Understanding the problem carefully}

If we only tried to fix parents to equal their children on the way up, a
parent larger than its children's total would have no legal move (we cannot
decrease). So first, descending, we \emph{raise} the children: give each
child at least the parent's value (bump the larger child, or either, up to
the parent). Now the children's total is $\ge$ the parent, so on the way
back up we can safely set each parent to the exact sum of its (possibly
raised) children.

\begin{ideabox}[Push Down, Then Fix Up]
\textbf{Descent:} at each node, compute the sum of its children's current
values; if that sum is $\ge$ the node, set the node to the sum, else push
the node's value down into (one of) its children. Recurse.
\textbf{Ascent:} after both children are finalised, set the node to the sum
of its children's values. Leaves are left untouched.
\end{ideabox}

\subsection*{Algorithm}

\begin{enumerate}
  \item \textsc{Fix}(\textit{node}):
  \begin{enumerate}
    \item If \textit{node} is null or a leaf, return.
    \item Let \textit{childSum} $=$ (left val if any) $+$ (right val if
          any). If \textit{childSum} $\ge$ \textit{node}.val, set
          \textit{node}.val $=$ \textit{childSum}; else copy
          \textit{node}.val into an existing child.
    \item Recurse into left and right.
    \item Set \textit{node}.val $=$ (left val if any) $+$ (right val if
          any).
  \end{enumerate}
\end{enumerate}

\subsection*{Dry run}

\dryrun{Tree: root $50$; children $7$ and $2$.}

\begin{itemize}
  \item Descent at $50$: childSum $=7+2=9<50$, so push $50$ down: set left
        child to $50$ (now $50,2$).
  \item Recurse into leaves $50$ and $2$: nothing to do.
  \item Ascent at root: set root $=50+2=52$.
  \item Final: root $52$, children $50,2$; $50+2=52$. Property holds.
\end{itemize}

\subsection*{C++ implementation}

\begin{lstlisting}
#include <bits/stdc++.h>
using namespace std;

struct TreeNode {
    int val;
    TreeNode* left;
    TreeNode* right;
    TreeNode(int v) : val(v), left(nullptr), right(nullptr) {}
};

void fix(TreeNode* node) {
    if (node == nullptr) return;
    if (node->left == nullptr && node->right == nullptr) return;

    int childSum = 0;
    if (node->left)  childSum += node->left->val;
    if (node->right) childSum += node->right->val;

    if (childSum >= node->val) {
        node->val = childSum;               // children already suffice
    } else {                                // push node value down
        if (node->left)  node->left->val  = node->val;
        else             node->right->val = node->val;
    }

    fix(node->left);
    fix(node->right);

    int total = 0;                          // fix up from actual children
    if (node->left)  total += node->left->val;
    if (node->right) total += node->right->val;
    node->val = total;
}

int main() {
    TreeNode* root = new TreeNode(50);
    root->left = new TreeNode(7);
    root->right = new TreeNode(2);

    fix(root);
    cout << root->val << " -> "
         << root->left->val << ", " << root->right->val << "\n";
    return 0;
}
\end{lstlisting}

\begin{complexitybox}
\textbf{Time Complexity.} $O(n)$ -- each node is handled once on the way
down and once on the way up. \textbf{Space Complexity.} $O(h)$ for the
recursion stack.
\end{complexitybox}

\begin{takeawaybox}
When a constraint restricts \emph{how} you may change values (here, only
increases), a single pass rarely suffices -- you pre-condition on the way
down so the fix-up on the way back is always legal. Down-then-up is a
recurring shape for tree mutations under monotone constraints.
\end{takeawaybox}

\section{All Nodes at Distance K from a Target}
\label{sec:trees-distance-k}

\problemstatement{Given the root, a target node, and an integer $K$, return
the values of all nodes that are exactly $K$ edges away from the target.
Distance may go upward (toward ancestors) as well as downward.}

\begin{patternbox}
The twist is \emph{``distance can go upward.''} A binary tree node has no
parent pointer, so radiating outward in all directions is impossible with
the raw structure. The fix: \textbf{build a parent map} so every node
effectively becomes a graph vertex with up-to-three neighbours (left
child, right child, parent), then run a plain \textbf{BFS} from the target.
\end{patternbox}

\subsection*{Understanding the problem carefully}

``All nodes at distance $K$'' is a BFS problem -- but BFS needs to move in
every direction, including up to the parent. So we first walk the tree
once recording each node's parent in a hash map. Now the tree behaves like
an undirected graph: from any node we can step to its left child, right
child, or parent. A breadth-first expansion from the target, stopping after
$K$ layers, lands exactly on the distance-$K$ nodes.

\begin{ideabox}[Turn the Tree into a Graph, Then BFS $K$ Layers]
Record parents with one DFS. Start BFS from the target with a
\textit{visited} set (so we never walk back the way we came). Expand
layer by layer; after $K$ layers the queue holds precisely the nodes at
distance $K$.
\end{ideabox}

\subsection*{Algorithm}

\begin{enumerate}
  \item DFS to fill \textit{parent}[node] for every node.
  \item BFS from target: mark it visited, distance $0$.
  \item Repeat $K$ times: expand the current frontier to all unvisited
        neighbours (left, right, parent), marking them visited.
  \item After $K$ expansions, the frontier is the answer.
\end{enumerate}

\subsection*{Dry run}

\dryrun{Tree: root $3$; left $5$ (children $6,2$); right $1$. Target $5$,
$K=2$.}

\begin{itemize}
  \item Parents: $5\!\to\!3$, $6\!\to\!5$, $2\!\to\!5$, $1\!\to\!3$.
  \item Layer $0$: $\{5\}$. Layer $1$: neighbours of $5$ = $6,2,3$
        (children and parent).
  \item Layer $2$: from $6$ nothing new; from $2$ nothing new; from $3$
        neighbours are $5$ (visited) and $1$ (new). Frontier $=\{1\}$.
  \item Nodes at distance $2$: $1$ (and, in the classic example, $7,4$ if
        present under $1$).
\end{itemize}

\subsection*{C++ implementation}

\begin{lstlisting}
#include <bits/stdc++.h>
using namespace std;

struct TreeNode {
    int val;
    TreeNode* left;
    TreeNode* right;
    TreeNode(int v) : val(v), left(nullptr), right(nullptr) {}
};

void markParents(TreeNode* node, unordered_map<TreeNode*, TreeNode*>& par) {
    if (node == nullptr) return;
    if (node->left)  { par[node->left]  = node; markParents(node->left,  par); }
    if (node->right) { par[node->right] = node; markParents(node->right, par); }
}

vector<int> distanceK(TreeNode* root, TreeNode* target, int K) {
    unordered_map<TreeNode*, TreeNode*> par;
    markParents(root, par);

    unordered_set<TreeNode*> visited;
    queue<TreeNode*> q;
    q.push(target);
    visited.insert(target);

    int dist = 0;
    while (!q.empty()) {
        if (dist == K) break;
        int s = q.size();
        for (int i = 0; i < s; i++) {
            TreeNode* node = q.front(); q.pop();
            for (TreeNode* nb : {node->left, node->right, par.count(node) ? par[node] : nullptr}) {
                if (nb && !visited.count(nb)) { visited.insert(nb); q.push(nb); }
            }
        }
        dist++;
    }

    vector<int> ans;
    while (!q.empty()) { ans.push_back(q.front()->val); q.pop(); }
    return ans;
}

int main() {
    TreeNode* root = new TreeNode(3);
    root->left = new TreeNode(5);
    root->right = new TreeNode(1);
    root->left->left  = new TreeNode(6);
    root->left->right = new TreeNode(2);

    vector<int> res = distanceK(root, root->left, 2);
    sort(res.begin(), res.end());
    for (int x : res) cout << x << " ";
    cout << "\n";
    return 0;
}
\end{lstlisting}

\begin{complexitybox}
\textbf{Time Complexity.} $O(n)$ -- one DFS to record parents and a BFS
visiting each node at most once. \textbf{Space Complexity.} $O(n)$ for the
parent map, visited set, and queue.
\end{complexitybox}

\begin{takeawaybox}
``Distance in all directions'' on a tree is a graph-BFS problem in
disguise; the enabling move is manufacturing parent pointers so upward
travel becomes possible. This parent-map trick recurs in the next section's
burning-tree problem.
\end{takeawaybox}

\section{Minimum Time to Burn the Tree}
\label{sec:trees-burn}

\problemstatement{A fire starts at a given target node. Each second, the
fire spreads from every burning node to its adjacent nodes (children and
parent). Return the minimum number of seconds needed to burn the entire
tree.}

\begin{patternbox}
This is the ``distance in all directions'' idea taken to its conclusion:
the answer is the \textbf{maximum distance} from the target to \emph{any}
node, because the last node to catch fire is the farthest one. Same
recipe -- \textbf{parent map plus BFS from the target} -- but instead of
stopping at layer $K$, count how many layers it takes to exhaust the tree.
\end{patternbox}

\subsection*{Understanding the problem carefully}

Fire spreads one edge per second in every direction. A node at graph
distance $d$ from the source catches fire at second $d$. The whole tree is
ablaze once the farthest node burns, so the time equals the eccentricity
of the target: the greatest distance to any node. A BFS from the target,
counting completed layers until the queue empties, measures exactly this.

\begin{ideabox}[Answer = Number of BFS Layers to Empty the Tree]
Build the parent map, BFS from the target with a visited set. Each fully
processed frontier is one second of spread. The number of layers after the
\emph{initial} one (the source itself, at time $0$) is the burn time --
i.e.\ increment a timer each time you expand a non-empty next frontier.
\end{ideabox}

\subsection*{Algorithm}

\begin{enumerate}
  \item DFS to fill \textit{parent}[node].
  \item BFS from target; \textit{time} $=-1$ (so the first layer, the
        source, corresponds to time $0$).
  \item Each layer: mark all unvisited neighbours; if the layer expanded
        anything, increment \textit{time}. Actually, increment
        \textit{time} once per processed frontier.
  \item Return \textit{time}.
\end{enumerate}

(Concretely: initialise \textit{time} $=0$, and increment it only when a
frontier produces at least one newly burning node.)

\subsection*{Dry run}

\dryrun{Tree: root $1$; left $2$ (left child $4$); right $3$. Start fire at
$4$.}

\begin{itemize}
  \item $t=0$: $\{4\}$ burning.
  \item $t=1$: neighbour of $4$ is $2$. $\{2\}$.
  \item $t=2$: neighbours of $2$ are $1$ (and $4$ burnt). $\{1\}$.
  \item $t=3$: neighbour of $1$ is $3$. $\{3\}$.
  \item Tree fully burnt at $t=3$.
\end{itemize}

\subsection*{C++ implementation}

\begin{lstlisting}
#include <bits/stdc++.h>
using namespace std;

struct TreeNode {
    int val;
    TreeNode* left;
    TreeNode* right;
    TreeNode(int v) : val(v), left(nullptr), right(nullptr) {}
};

void markParents(TreeNode* node, unordered_map<TreeNode*, TreeNode*>& par) {
    if (node == nullptr) return;
    if (node->left)  { par[node->left]  = node; markParents(node->left,  par); }
    if (node->right) { par[node->right] = node; markParents(node->right, par); }
}

int timeToBurn(TreeNode* root, TreeNode* target) {
    unordered_map<TreeNode*, TreeNode*> par;
    markParents(root, par);

    unordered_set<TreeNode*> visited;
    queue<TreeNode*> q;
    q.push(target);
    visited.insert(target);

    int time = 0;
    while (!q.empty()) {
        int s = q.size();
        bool spread = false;
        for (int i = 0; i < s; i++) {
            TreeNode* node = q.front(); q.pop();
            for (TreeNode* nb : {node->left, node->right, par.count(node) ? par[node] : nullptr}) {
                if (nb && !visited.count(nb)) {
                    visited.insert(nb);
                    q.push(nb);
                    spread = true;
                }
            }
        }
        if (spread) time++;
    }
    return time;
}

int main() {
    TreeNode* root = new TreeNode(1);
    root->left = new TreeNode(2);
    root->right = new TreeNode(3);
    root->left->left = new TreeNode(4);

    cout << "Burn time: " << timeToBurn(root, root->left->left) << "\n";
    return 0;
}
\end{lstlisting}

\begin{complexitybox}
\textbf{Time Complexity.} $O(n)$ -- one DFS plus one BFS, each visiting
every node once. \textbf{Space Complexity.} $O(n)$ for the parent map,
visited set, and queue.
\end{complexitybox}

\begin{takeawaybox}
Burning the tree is ``nodes at distance $K$'' with $K$ pushed to the
maximum: the answer is the number of BFS layers to consume the whole tree.
Once you see fire-spread problems as level-counting BFS on the
parent-augmented graph, they become routine.
\end{takeawaybox}

\section{Count Nodes in a Complete Binary Tree}
\label{sec:trees-count-complete}

\problemstatement{Given the root of a \emph{complete} binary tree (every
level full except possibly the last, which is filled left to right), count
its nodes in better than $O(n)$ time.}

\begin{patternbox}
A plain traversal counts in $O(n)$, but the word \emph{complete} is an
invitation to do better. The key fact: if a subtree's leftmost depth equals
its rightmost depth, it is a \textbf{perfect} tree and its node count is
$2^{\text{height}}-1$ by formula -- no traversal needed. Only along the
one ``ragged'' edge do we recurse, giving $O(\log^2 n)$.
\end{patternbox}

\subsection*{Understanding the problem carefully}

Walk down the far left of a subtree to get its left height; walk down the
far right to get its right height. If they are equal, every level is full,
so the subtree is perfect and has exactly $2^h-1$ nodes -- return that
instantly. If they differ, the last level is only partially filled, so
recurse into both children and add $1$ for the current node. Crucially, at
most one child per level is ragged; the other is perfect and resolved by
formula, so the recursion depth is $O(\log n)$ and each level costs
$O(\log n)$ to measure the two heights.

\begin{ideabox}[Perfect Subtrees Close by Formula]
\textsc{Count}(node): if null return $0$. Measure left height $l$
(all-left walk) and right height $r$ (all-right walk). If $l=r$, return
$2^{l}-1$. Otherwise return $1 + \textsc{Count}(\text{left}) +
\textsc{Count}(\text{right})$.
\end{ideabox}

\subsection*{Algorithm}

\begin{enumerate}
  \item \textsc{leftHeight}(node): count nodes going only left until null.
  \item \textsc{rightHeight}(node): count nodes going only right until
        null.
  \item \textsc{Count}(node): if null return $0$; if leftHeight $=$
        rightHeight $=h$ return $2^h-1$; else return $1+$
        \textsc{Count}(left)$+$\textsc{Count}(right).
\end{enumerate}

\subsection*{Dry run}

\dryrun{Complete tree: root $1$; children $2,3$; $2$'s children $4,5$;
$3$'s left child $6$.}

\begin{itemize}
  \item At $1$: left height $=3$ ($1\!\to\!2\!\to\!4$), right height $=2$
        ($1\!\to\!3$). Unequal $\Rightarrow$ recurse.
  \item Left subtree at $2$: left height $=$ right height $=2$; perfect,
        $2^2-1=3$ nodes.
  \item Right subtree at $3$: left height $=2$, right height $=1$; recurse
        $\to 1+\textsc{Count}(6)+0 = 1+1 = 2$.
  \item Total $=1+3+2=6$.
\end{itemize}

\subsection*{C++ implementation}

\begin{lstlisting}
#include <bits/stdc++.h>
using namespace std;

struct TreeNode {
    int val;
    TreeNode* left;
    TreeNode* right;
    TreeNode(int v) : val(v), left(nullptr), right(nullptr) {}
};

int leftHeight(TreeNode* n) {
    int h = 0;
    while (n) { h++; n = n->left; }
    return h;
}
int rightHeight(TreeNode* n) {
    int h = 0;
    while (n) { h++; n = n->right; }
    return h;
}

int countNodes(TreeNode* root) {
    if (root == nullptr) return 0;
    int lh = leftHeight(root), rh = rightHeight(root);
    if (lh == rh) return (1 << lh) - 1;          // perfect subtree
    return 1 + countNodes(root->left) + countNodes(root->right);
}

int main() {
    TreeNode* root = new TreeNode(1);
    root->left = new TreeNode(2);
    root->right = new TreeNode(3);
    root->left->left = new TreeNode(4);
    root->left->right = new TreeNode(5);
    root->right->left = new TreeNode(6);

    cout << "Node count: " << countNodes(root) << "\n";
    return 0;
}
\end{lstlisting}

\begin{complexitybox}
\textbf{Time Complexity.} $O(\log^2 n)$ -- the recursion has $O(\log n)$
depth, and each level spends $O(\log n)$ measuring the two heights.
\textbf{Space Complexity.} $O(\log n)$ for the recursion stack.
\end{complexitybox}

\begin{takeawaybox}
When a structural guarantee (``complete'') is handed to you, look for a
closed-form shortcut before defaulting to a full traversal. Resolving
perfect subtrees by the $2^h-1$ formula and recursing only along the ragged
boundary is the canonical sub-linear tree count.
\end{takeawaybox}

\section{Construct a Binary Tree from Preorder and Inorder}
\label{sec:trees-construct-pre-in}

\problemstatement{Given the preorder and inorder traversals of a binary
tree with unique values, reconstruct and return the tree.}

\begin{patternbox}
\emph{``Rebuild a tree from two traversals''} is a divide-and-conquer:
\textbf{preorder's first element is always the root}, and locating that
root inside the inorder array \textbf{splits inorder into the left subtree
and the right subtree}. Recurse on the two halves. A hash map from value to
inorder index turns the ``locate the root'' step from $O(n)$ into $O(1)$.
\end{patternbox}

\subsection*{Understanding the problem carefully}

Preorder visits root, then all of left, then all of right -- so its first
value is the root. Inorder visits left, root, right -- so once we know the
root's value, everything to its left in the inorder array is the left
subtree and everything to its right is the right subtree. The count of
left-subtree nodes (from inorder) tells us exactly how many of the
following preorder elements belong to the left subtree, letting us slice
preorder too. Recursing reconstructs both subtrees.

\begin{ideabox}[Preorder Gives the Root; Inorder Gives the Split]
Take the next preorder value as the current root. Find it in inorder at
position $p$: the $p$ elements left of it form the left subtree, the rest
form the right subtree. Consume preorder left-to-right (root, then build
left, then build right). An index map makes finding $p$ instant.
\end{ideabox}

\subsection*{Algorithm}

\begin{enumerate}
  \item Build \textit{pos}: value $\to$ its index in inorder.
  \item \textsc{Build}(\textit{inLo}, \textit{inHi}) using a shared
        preorder cursor \textit{preIdx} $=0$:
  \begin{enumerate}
    \item If \textit{inLo} $>$ \textit{inHi}, return null.
    \item \textit{rootVal} $=$ preorder[\textit{preIdx}++]; make a node.
    \item $p=\textit{pos}[\textit{rootVal}]$.
    \item node.left $=$ \textsc{Build}(\textit{inLo}, $p-1$); node.right
          $=$ \textsc{Build}($p+1$, \textit{inHi}).
    \item Return node.
  \end{enumerate}
\end{enumerate}

\subsection*{Dry run}

\dryrun{preorder $=[3,9,20,15,7]$, inorder $=[9,3,15,20,7]$.}

\begin{itemize}
  \item Root $=3$ (preorder[0]); in inorder at index $1$: left $=[9]$,
        right $=[15,20,7]$.
  \item Left: root $=9$; inorder slice $[9]$ has no children -- leaf.
  \item Right: root $=20$ (next preorder); in inorder index $3$: left
        $=[15]$, right $=[7]$; both leaves.
  \item Reconstructed: $3$ with left $9$ and right $20(15,7)$.
\end{itemize}

\subsection*{C++ implementation}

\begin{lstlisting}
#include <bits/stdc++.h>
using namespace std;

struct TreeNode {
    int val;
    TreeNode* left;
    TreeNode* right;
    TreeNode(int v) : val(v), left(nullptr), right(nullptr) {}
};

int preIdx;

TreeNode* build(vector<int>& preorder, int inLo, int inHi,
                unordered_map<int,int>& pos) {
    if (inLo > inHi) return nullptr;
    int rootVal = preorder[preIdx++];
    TreeNode* node = new TreeNode(rootVal);
    int p = pos[rootVal];
    node->left  = build(preorder, inLo, p - 1, pos);
    node->right = build(preorder, p + 1, inHi, pos);
    return node;
}

TreeNode* buildTree(vector<int>& preorder, vector<int>& inorder) {
    unordered_map<int,int> pos;
    for (int i = 0; i < (int)inorder.size(); i++) pos[inorder[i]] = i;
    preIdx = 0;
    return build(preorder, 0, inorder.size() - 1, pos);
}

void inorderPrint(TreeNode* n) {
    if (!n) return;
    inorderPrint(n->left);
    cout << n->val << " ";
    inorderPrint(n->right);
}

int main() {
    vector<int> preorder = {3, 9, 20, 15, 7};
    vector<int> inorder  = {9, 3, 15, 20, 7};
    TreeNode* root = buildTree(preorder, inorder);
    inorderPrint(root);            // should reprint the inorder
    cout << "\n";
    return 0;
}
\end{lstlisting}

\begin{complexitybox}
\textbf{Time Complexity.} $O(n)$ -- each node is created once and the
inorder position is an $O(1)$ map lookup. \textbf{Space Complexity.}
$O(n)$ for the map, plus $O(h)$ recursion.
\end{complexitybox}

\begin{takeawaybox}
Reconstruction hinges on two facts: one traversal names the root, another
splits the rest into left and right. The shared preorder cursor consumed
left-to-right is the clean way to avoid slicing arrays at every recursive
call.
\end{takeawaybox}

\section{Construct a Binary Tree from Postorder and Inorder}
\label{sec:trees-construct-post-in}

\problemstatement{Given the postorder and inorder traversals of a binary
tree with unique values, reconstruct and return the tree.}

\begin{patternbox}
This is the previous section mirrored. Postorder visits left, right,
\emph{root} -- so the root is the \textbf{last} element, not the first. Read
postorder \emph{right to left} and it behaves like preorder with the
children swapped: root, then right subtree, then left subtree. So build the
\textbf{right subtree before the left}, consuming postorder from the back.
\end{patternbox}

\subsection*{Understanding the problem carefully}

Because postorder ends with the root, a cursor moving backward through
postorder yields roots in the order root, right-subtree-root,
\ldots. Inorder still splits into left and right around the root's
position. The one change from the preorder version: since we now encounter
the right subtree's root immediately after the current root (reading
backward), we must recurse \emph{right first}, then left, so the backward
cursor lines up with the subtrees correctly.

\begin{ideabox}[Postorder from the Back Names Roots; Build Right, Then Left]
Take the current root as postorder[\textit{postIdx}] and decrement
\textit{postIdx}. Split inorder at the root's position. Recurse into the
\emph{right} subtree first (node.right), then the \emph{left}
(node.left) -- the reverse of the preorder construction.
\end{ideabox}

\subsection*{Algorithm}

\begin{enumerate}
  \item Build \textit{pos}: value $\to$ inorder index. Set
        \textit{postIdx} $=n-1$.
  \item \textsc{Build}(\textit{inLo}, \textit{inHi}):
  \begin{enumerate}
    \item If \textit{inLo} $>$ \textit{inHi}, return null.
    \item \textit{rootVal} $=$ postorder[\textit{postIdx}--]; make node.
    \item $p=\textit{pos}[\textit{rootVal}]$.
    \item node.right $=$ \textsc{Build}($p+1$, \textit{inHi}); node.left
          $=$ \textsc{Build}(\textit{inLo}, $p-1$).
    \item Return node.
  \end{enumerate}
\end{enumerate}

\subsection*{Dry run}

\dryrun{postorder $=[9,15,7,20,3]$, inorder $=[9,3,15,20,7]$.}

\begin{itemize}
  \item Root $=3$ (postorder last); inorder index $1$: left $=[9]$, right
        $=[15,20,7]$.
  \item Build right first: next root (moving back) $=20$; inorder index
        $3$: right $=[7]$ (leaf), left $=[15]$ (leaf).
  \item Build left: next root $=9$ -- leaf.
  \item Reconstructed tree matches the original.
\end{itemize}

\subsection*{C++ implementation}

\begin{lstlisting}
#include <bits/stdc++.h>
using namespace std;

struct TreeNode {
    int val;
    TreeNode* left;
    TreeNode* right;
    TreeNode(int v) : val(v), left(nullptr), right(nullptr) {}
};

int postIdx;

TreeNode* build(vector<int>& postorder, int inLo, int inHi,
                unordered_map<int,int>& pos) {
    if (inLo > inHi) return nullptr;
    int rootVal = postorder[postIdx--];
    TreeNode* node = new TreeNode(rootVal);
    int p = pos[rootVal];
    node->right = build(postorder, p + 1, inHi, pos);  // right first
    node->left  = build(postorder, inLo, p - 1, pos);
    return node;
}

TreeNode* buildTree(vector<int>& inorder, vector<int>& postorder) {
    unordered_map<int,int> pos;
    for (int i = 0; i < (int)inorder.size(); i++) pos[inorder[i]] = i;
    postIdx = postorder.size() - 1;
    return build(postorder, 0, inorder.size() - 1, pos);
}

void inorderPrint(TreeNode* n) {
    if (!n) return;
    inorderPrint(n->left);
    cout << n->val << " ";
    inorderPrint(n->right);
}

int main() {
    vector<int> inorder   = {9, 3, 15, 20, 7};
    vector<int> postorder = {9, 15, 7, 20, 3};
    TreeNode* root = buildTree(inorder, postorder);
    inorderPrint(root);
    cout << "\n";
    return 0;
}
\end{lstlisting}

\begin{complexitybox}
\textbf{Time Complexity.} $O(n)$ with the inorder index map. \textbf{Space
Complexity.} $O(n)$ for the map plus $O(h)$ recursion.
\end{complexitybox}

\begin{takeawaybox}
Postorder+inorder is preorder+inorder read backward with the child order
flipped. Seeing the symmetry means you derive one from the other instead of
memorising two separate algorithms. Note: preorder+postorder alone cannot
uniquely rebuild a tree -- inorder is what fixes the left/right split.
\end{takeawaybox}

\section{Serialize and Deserialize a Binary Tree}
\label{sec:trees-serialize}

\problemstatement{Design two functions: \emph{serialize}, which encodes a
binary tree into a single string, and \emph{deserialize}, which rebuilds
the exact tree from that string. Values may repeat and the tree may be
arbitrary.}

\begin{patternbox}
Unlike reconstruction from two traversals, here we control the format --
so we choose one that is \emph{self-describing} by \textbf{recording null
children explicitly}. With null markers, a \emph{single} preorder walk
carries all the structure needed to rebuild the tree unambiguously. The
design freedom is the whole point.
\end{patternbox}

\subsection*{Understanding the problem carefully}

Two traversals are needed to rebuild a tree only because a normal traversal
hides where subtrees end. If we instead emit a special token (say \verb|#|)
whenever we hit a null child, the string records the tree's exact shape.
Reading that string back in the same preorder order, a null token means
``this child is empty,'' so deserialization mirrors serialization step for
step -- read a token, build a node (or return null), then recursively build
its left and right children.

\begin{ideabox}[Preorder with Explicit Null Markers Is Self-Sufficient]
\textbf{Serialize:} preorder traversal; output each value, and output a
sentinel \verb|#| for every null child, separated by commas.
\textbf{Deserialize:} read tokens in order; \verb|#| yields null,
otherwise create a node and recursively deserialize its left then right
child. The null markers tell the reader exactly where each subtree stops.
\end{ideabox}

\subsection*{Algorithm}

\begin{enumerate}
  \item \textsc{Serialize}(node, out): if node is null, append
        \verb|#,|; else append \verb|node.val,| then
        \textsc{Serialize}(left) and \textsc{Serialize}(right).
  \item \textsc{Deserialize}: read the next token; if \verb|#| return
        null; else make a node, set node.left $=$ \textsc{Deserialize},
        node.right $=$ \textsc{Deserialize}; return node.
\end{enumerate}

\subsection*{Dry run}

\dryrun{Tree: root $1$; left $2$; right $3$ (with children $4,5$).}

\begin{itemize}
  \item Serialize: \verb|1,2,#,#,3,4,#,#,5,#,#,| -- $2$ is a leaf (two
        \verb|#|), then $3$, then its children $4,5$.
  \item Deserialize reads $1$ (node), builds left from $2$ (then two
        \verb|#| $\to$ leaf), builds right from $3$, whose children come
        from $4$ and $5$.
  \item The rebuilt tree is identical to the original.
\end{itemize}

\subsection*{C++ implementation}

\begin{lstlisting}
#include <bits/stdc++.h>
using namespace std;

struct TreeNode {
    int val;
    TreeNode* left;
    TreeNode* right;
    TreeNode(int v) : val(v), left(nullptr), right(nullptr) {}
};

void serialize(TreeNode* node, string& out) {
    if (node == nullptr) { out += "#,"; return; }
    out += to_string(node->val) + ",";
    serialize(node->left, out);
    serialize(node->right, out);
}

TreeNode* deserialize(istringstream& in) {
    string token;
    getline(in, token, ',');
    if (token == "#") return nullptr;
    TreeNode* node = new TreeNode(stoi(token));
    node->left  = deserialize(in);
    node->right = deserialize(in);
    return node;
}

void inorderPrint(TreeNode* n) {
    if (!n) return;
    inorderPrint(n->left);
    cout << n->val << " ";
    inorderPrint(n->right);
}

int main() {
    TreeNode* root = new TreeNode(1);
    root->left = new TreeNode(2);
    root->right = new TreeNode(3);
    root->right->left = new TreeNode(4);
    root->right->right = new TreeNode(5);

    string s;
    serialize(root, s);
    cout << "Serialized: " << s << "\n";

    istringstream in(s);
    TreeNode* rebuilt = deserialize(in);
    cout << "Inorder of rebuilt: ";
    inorderPrint(rebuilt);
    cout << "\n";
    return 0;
}
\end{lstlisting}

\begin{complexitybox}
\textbf{Time Complexity.} $O(n)$ for both serialize and deserialize --
each node and each null marker is written and read once. \textbf{Space
Complexity.} $O(n)$ for the string, plus $O(h)$ recursion.
\end{complexitybox}

\begin{takeawaybox}
When you own the encoding, adding explicit null markers makes a single
traversal invertible -- no second traversal required. This ``sentinel for
absence'' idea is the general trick for making any recursive structure
round-trip through a flat string.
\end{takeawaybox}

\section{Morris Inorder Traversal}
\label{sec:trees-morris-inorder}

\problemstatement{Return the inorder traversal of a binary tree using
$O(1)$ extra space -- no recursion and no explicit stack.}

\begin{patternbox}
\emph{``$O(1)$ space,'' ``no stack,'' ``no recursion''} points to
\textbf{Morris traversal}, which recovers the ``go back up to the
ancestor'' capability -- normally provided by a stack -- by temporarily
rewiring the tree's own unused null pointers into \textbf{threads} back to
the ancestor, then undoing them.
\end{patternbox}

\subsection*{Understanding the problem carefully}

The only reason inorder needs a stack is to remember, after finishing a
left subtree, which node to return to. Morris's insight: the node we must
return to is exactly the \emph{inorder predecessor}'s successor -- and the
inorder predecessor of a node (the rightmost node of its left subtree) has
a null right pointer we can borrow. We point that null back at the current
node, creating a temporary ``thread.'' After the left subtree is walked, we
arrive back via the thread, record the node, remove the thread (restoring
the tree), and move right.

\begin{ideabox}[Thread the Predecessor's Null Right Pointer Back to You]
At each node with a left child, find its inorder predecessor (go left once,
then all the way right). If the predecessor's right is null, set it to the
current node (make the thread) and move left. If it is already the current
node, the left subtree is done: remove the thread, record the node, and
move right. Nodes without a left child are recorded immediately and you
move right.
\end{ideabox}

\subsection*{Algorithm}

\begin{enumerate}
  \item \textit{curr} $=$ root.
  \item While \textit{curr} is non-null:
  \begin{enumerate}
    \item If \textit{curr} has no left child: record \textit{curr}; move
          \textit{curr} $=$ \textit{curr}.right.
    \item Else find \textit{pred} $=$ rightmost node of
          \textit{curr}.left (not passing through an existing thread):
    \begin{enumerate}
      \item If \textit{pred}.right is null: set it to \textit{curr};
            move \textit{curr} $=$ \textit{curr}.left.
      \item Else (\textit{pred}.right $=$ \textit{curr}): reset
            \textit{pred}.right to null; record \textit{curr}; move
            \textit{curr} $=$ \textit{curr}.right.
    \end{enumerate}
  \end{enumerate}
\end{enumerate}

\subsection*{Dry run}

\dryrun{Tree: root $2$; left $1$; right $3$.}

\begin{itemize}
  \item \textit{curr}$=2$, has left $1$. Predecessor of $2$ is $1$;
        $1$.right null $\to$ thread $1\!\to\!2$; move left to $1$.
  \item \textit{curr}$=1$, no left $\to$ record $1$; move right via
        thread to $2$.
  \item \textit{curr}$=2$, has left; predecessor $1$.right $=2$ (thread)
        $\to$ remove thread, record $2$, move right to $3$.
  \item \textit{curr}$=3$, no left $\to$ record $3$.
  \item Inorder: $1,2,3$.
\end{itemize}

\subsection*{C++ implementation}

\begin{lstlisting}
#include <bits/stdc++.h>
using namespace std;

struct TreeNode {
    int val;
    TreeNode* left;
    TreeNode* right;
    TreeNode(int v) : val(v), left(nullptr), right(nullptr) {}
};

vector<int> morrisInorder(TreeNode* root) {
    vector<int> out;
    TreeNode* curr = root;
    while (curr != nullptr) {
        if (curr->left == nullptr) {
            out.push_back(curr->val);        // no left: record, go right
            curr = curr->right;
        } else {
            TreeNode* pred = curr->left;
            while (pred->right != nullptr && pred->right != curr)
                pred = pred->right;          // find inorder predecessor
            if (pred->right == nullptr) {
                pred->right = curr;          // create thread
                curr = curr->left;
            } else {
                pred->right = nullptr;       // remove thread
                out.push_back(curr->val);    // left subtree done: record
                curr = curr->right;
            }
        }
    }
    return out;
}

int main() {
    TreeNode* root = new TreeNode(2);
    root->left = new TreeNode(1);
    root->right = new TreeNode(3);

    for (int x : morrisInorder(root)) cout << x << " ";
    cout << "\n";
    return 0;
}
\end{lstlisting}

\begin{complexitybox}
\textbf{Time Complexity.} $O(n)$ -- although finding predecessors seems to
add work, each edge is traversed at most a constant number of times, so
the total remains linear. \textbf{Space Complexity.} $O(1)$ -- no stack or
recursion; only a few pointers, with the tree restored to its original
shape at the end.
\end{complexitybox}

\begin{takeawaybox}
Morris trades the stack for temporary threads woven from the tree's own
null pointers, achieving true $O(1)$ space. The discipline that keeps it
correct is always removing a thread the second time you meet it, so the
tree is left exactly as it began.
\end{takeawaybox}

\section{Morris Preorder Traversal}
\label{sec:trees-morris-preorder}

\problemstatement{Return the preorder traversal of a binary tree in $O(1)$
extra space -- no recursion and no explicit stack.}

\begin{patternbox}
Morris preorder is Morris inorder with the ``record'' step moved. Recall
from Section~\ref{sec:trees-rec-traversals} that preorder and inorder are
the same walk with the node recorded at a different moment. In the threaded
walk, that means recording the node \emph{when you create the thread}
(first visit) rather than when you remove it.
\end{patternbox}

\subsection*{Understanding the problem carefully}

The Morris machinery -- thread the predecessor's null right pointer, walk
left, later return and unthread -- visits each node twice when it has a left
child: once on the way in (creating the thread) and once on the way back
(removing it). Inorder records on the second visit; preorder must record on
the \emph{first}. So we output the node's value at the moment we create the
thread (and immediately for nodes with no left child), which places the
node before its subtrees -- exactly preorder.

\begin{ideabox}[Same Threads, Record on the First Visit]
Identical to Morris inorder, with one line moved: when you create a thread
(first time you reach a node with a left child), record the node
\emph{then}, before moving left. Nodes without a left child are still
recorded immediately. When you later remove a thread, do \emph{not} record
again.
\end{ideabox}

\subsection*{Algorithm}

\begin{enumerate}
  \item \textit{curr} $=$ root.
  \item While \textit{curr} is non-null:
  \begin{enumerate}
    \item If \textit{curr} has no left child: record \textit{curr}; move
          right.
    \item Else find predecessor \textit{pred} (rightmost of the left
          subtree, not through a thread):
    \begin{enumerate}
      \item If \textit{pred}.right is null: \textbf{record
            \textit{curr}} (first visit); set \textit{pred}.right $=$
            \textit{curr}; move left.
      \item Else: reset \textit{pred}.right to null; move right (no
            recording).
    \end{enumerate}
  \end{enumerate}
\end{enumerate}

\subsection*{Dry run}

\dryrun{Tree: root $1$; left $2$; right $3$.}

\begin{itemize}
  \item \textit{curr}$=1$, has left. Predecessor $2$.right null $\to$
        record $1$; thread $2\!\to\!1$; move left to $2$.
  \item \textit{curr}$=2$, no left $\to$ record $2$; move right via
        thread to $1$.
  \item \textit{curr}$=1$, predecessor $2$.right $=1$ (thread) $\to$
        remove thread; move right to $3$ (no record).
  \item \textit{curr}$=3$, no left $\to$ record $3$.
  \item Preorder: $1,2,3$.
\end{itemize}

\subsection*{C++ implementation}

\begin{lstlisting}
#include <bits/stdc++.h>
using namespace std;

struct TreeNode {
    int val;
    TreeNode* left;
    TreeNode* right;
    TreeNode(int v) : val(v), left(nullptr), right(nullptr) {}
};

vector<int> morrisPreorder(TreeNode* root) {
    vector<int> out;
    TreeNode* curr = root;
    while (curr != nullptr) {
        if (curr->left == nullptr) {
            out.push_back(curr->val);
            curr = curr->right;
        } else {
            TreeNode* pred = curr->left;
            while (pred->right != nullptr && pred->right != curr)
                pred = pred->right;
            if (pred->right == nullptr) {
                out.push_back(curr->val);    // record on FIRST visit
                pred->right = curr;
                curr = curr->left;
            } else {
                pred->right = nullptr;       // second visit: just unthread
                curr = curr->right;
            }
        }
    }
    return out;
}

int main() {
    TreeNode* root = new TreeNode(1);
    root->left = new TreeNode(2);
    root->right = new TreeNode(3);

    for (int x : morrisPreorder(root)) cout << x << " ";
    cout << "\n";
    return 0;
}
\end{lstlisting}

\begin{complexitybox}
\textbf{Time Complexity.} $O(n)$ -- each edge is walked a constant number
of times. \textbf{Space Complexity.} $O(1)$ -- threads only, tree restored
at the end.
\end{complexitybox}

\begin{takeawaybox}
Once Morris inorder is understood, preorder is a one-line change -- record
on the first visit instead of the second -- reaffirming the chapter's
opening theme that the three depth-first orders differ only in \emph{when}
the node is recorded, threading or not.
\end{takeawaybox}

\section{Flatten a Binary Tree to a Linked List}
\label{sec:trees-flatten}

\problemstatement{Flatten a binary tree in place so that it becomes a
``linked list'': every node's left child is null and its right child points
to the next node in \emph{preorder}. Use the existing right pointers as the
list's next pointers.}

\begin{patternbox}
The target order is \emph{preorder}, so the flattened list is node, then
all of the (flattened) left subtree, then all of the (flattened) right
subtree. The elegant $O(1)$-space solution is a \textbf{reverse-preorder}
(right, left, node) walk that stitches each node's right pointer to the
previously visited node -- building the list backwards.
\end{patternbox}

\subsection*{Understanding the problem carefully}

In the final list, after any node come all its left-subtree nodes, then its
right-subtree nodes. If we process nodes in the \emph{reverse} of preorder
-- that is right subtree, then left subtree, then the node -- then at the
moment we handle a node, the entire portion of the list that should come
\emph{after} it has already been assembled and is pointed to by a running
\textit{prev} pointer. So we set node.right $=$ \textit{prev}, node.left
$=$ null, and update \textit{prev} $=$ node.

\begin{ideabox}[Build the List Backwards with a Running \textit{prev}]
Traverse in right, left, node order. Keep \textit{prev}, the head of the
already-flattened suffix. At each node: set its right pointer to
\textit{prev}, its left pointer to null, then set \textit{prev} to this
node. When the traversal finishes, \textit{prev} is the root of the fully
flattened list.
\end{ideabox}

\subsection*{Algorithm}

\begin{enumerate}
  \item \textit{prev} $=$ null.
  \item \textsc{Flatten}(node):
  \begin{enumerate}
    \item If node is null, return.
    \item \textsc{Flatten}(node.right).
    \item \textsc{Flatten}(node.left).
    \item node.right $=$ \textit{prev}; node.left $=$ null; \textit{prev}
          $=$ node.
  \end{enumerate}
\end{enumerate}

\subsection*{Dry run}

\dryrun{Tree: root $1$; left $2$ (children $3,4$); right $5$ (right child
$6$). Preorder: $1,2,3,4,5,6$.}

\begin{itemize}
  \item Reverse-preorder processes $6$ first: $6$.right$=$null,
        \textit{prev}$=6$.
  \item Then $5$: $5$.right$=6$, \textit{prev}$=5$. Then $4$:
        $4$.right$=5$, \textit{prev}$=4$. Then $3$: $3$.right$=4$. Then
        $2$: $2$.right$=3$, left null. Then $1$: $1$.right$=2$, left null.
  \item Result chain via right pointers: $1\!\to\!2\!\to\!3\!\to\!4\!\to
        \!5\!\to\!6$, all left pointers null.
\end{itemize}

\subsection*{C++ implementation}

\begin{lstlisting}
#include <bits/stdc++.h>
using namespace std;

struct TreeNode {
    int val;
    TreeNode* left;
    TreeNode* right;
    TreeNode(int v) : val(v), left(nullptr), right(nullptr) {}
};

TreeNode* prevNode = nullptr;

void flatten(TreeNode* node) {
    if (node == nullptr) return;
    flatten(node->right);      // right first
    flatten(node->left);       // then left
    node->right = prevNode;    // stitch onto the already-built suffix
    node->left = nullptr;
    prevNode = node;
}

int main() {
    TreeNode* root = new TreeNode(1);
    root->left = new TreeNode(2);
    root->right = new TreeNode(5);
    root->left->left = new TreeNode(3);
    root->left->right = new TreeNode(4);
    root->right->right = new TreeNode(6);

    prevNode = nullptr;
    flatten(root);

    for (TreeNode* c = root; c != nullptr; c = c->right)
        cout << c->val << " ";
    cout << "\n";
    return 0;
}
\end{lstlisting}

\begin{complexitybox}
\textbf{Time Complexity.} $O(n)$ -- each node is visited once.
\textbf{Space Complexity.} $O(h)$ for the recursion stack; a Morris-style
variant achieves $O(1)$ by rewiring the rightmost node of each left subtree
to the right subtree, in the spirit of the previous two sections.
\end{complexitybox}

\begin{takeawaybox}
Processing nodes in reverse of the target order, while carrying a
\textit{prev} pointer to the already-built tail, is a recurring in-place
linking technique -- the same idea that reverses a linked list. Flattening
a tree is that trick applied along a preorder spine.
\end{takeawaybox}

